% Options for packages loaded elsewhere
\PassOptionsToPackage{unicode}{hyperref}
\PassOptionsToPackage{hyphens}{url}
\PassOptionsToPackage{dvipsnames,svgnames,x11names}{xcolor}
\documentclass[
]{article}
\usepackage{xcolor}
\usepackage[margin=1in]{geometry}
\usepackage{amsmath,amssymb}
\setcounter{secnumdepth}{-\maxdimen} % remove section numbering
\usepackage{iftex}
\ifPDFTeX
  \usepackage[T1]{fontenc}
  \usepackage[utf8]{inputenc}
  \usepackage{textcomp} % provide euro and other symbols
\else % if luatex or xetex
  \usepackage{unicode-math} % this also loads fontspec
  \defaultfontfeatures{Scale=MatchLowercase}
  \defaultfontfeatures[\rmfamily]{Ligatures=TeX,Scale=1}
\fi
\usepackage{lmodern}
\ifPDFTeX\else
  % xetex/luatex font selection
\fi
% Use upquote if available, for straight quotes in verbatim environments
\IfFileExists{upquote.sty}{\usepackage{upquote}}{}
\IfFileExists{microtype.sty}{% use microtype if available
  \usepackage[]{microtype}
  \UseMicrotypeSet[protrusion]{basicmath} % disable protrusion for tt fonts
}{}
\makeatletter
\@ifundefined{KOMAClassName}{% if non-KOMA class
  \IfFileExists{parskip.sty}{%
    \usepackage{parskip}
  }{% else
    \setlength{\parindent}{0pt}
    \setlength{\parskip}{6pt plus 2pt minus 1pt}}
}{% if KOMA class
  \KOMAoptions{parskip=half}}
\makeatother
\usepackage{longtable,booktabs,array}
\newcounter{none} % for unnumbered tables
\usepackage{calc} % for calculating minipage widths
% Correct order of tables after \paragraph or \subparagraph
\usepackage{etoolbox}
\makeatletter
\patchcmd\longtable{\par}{\if@noskipsec\mbox{}\fi\par}{}{}
\makeatother
% Allow footnotes in longtable head/foot
\IfFileExists{footnotehyper.sty}{\usepackage{footnotehyper}}{\usepackage{footnote}}
\makesavenoteenv{longtable}
\setlength{\emergencystretch}{3em} % prevent overfull lines
\providecommand{\tightlist}{%
  \setlength{\itemsep}{0pt}\setlength{\parskip}{0pt}}
\usepackage{bookmark}
\IfFileExists{xurl.sty}{\usepackage{xurl}}{} % add URL line breaks if available
\urlstyle{same}
\hypersetup{
  pdftitle={Told or Enforced: When In-Context Contracts Substitute for Runtime Enforcement in Agent Harnesses},
  pdfauthor={Dom Colligan},
  colorlinks=true,
  linkcolor={Maroon},
  filecolor={Maroon},
  citecolor={Blue},
  urlcolor={Blue},
  pdfcreator={LaTeX via pandoc}}

\title{Told or Enforced: When In-Context Contracts Substitute for
Runtime Enforcement in Agent Harnesses}
\author{Dom Colligan}
\date{}

\begin{document}
\maketitle

{
\setcounter{tocdepth}{3}
\tableofcontents
}
\section{Told or Enforced: When In-Context Contracts Substitute for
Runtime Enforcement in Agent
Harnesses}\label{told-or-enforced-when-in-context-contracts-substitute-for-runtime-enforcement-in-agent-harnesses}

Dom Colligan

\subsection{Abstract}\label{abstract}

For each constraint, a harness designer chooses: state the rule in the
agent's prompt (an in-context contract), or have the runtime make the
violation impossible (runtime enforcement). We cross the two levers in a
2x2 per governance mechanism of a frozen reference runtime, separating
the effect of telling, the effect of enforcing, and the marginal value
of enforcement given the agent was told. In a small, pre-registered,
single-runtime case study we find a capability-moderated pattern on the
mutation gate (the runtime's user-commit boundary): a model that can
self-enforce the boundary respects it with the runtime off, while a
model that cannot needs the runtime to do the whole job.

The evidence is a four-model ladder (MiniMax-M3 and Llama-3.3-70B as the
capable rungs; Qwen3.5-9B and Qwen2.5-7B as the weak rungs). Told the
rule with the runtime off, both capable models complied on every
replicate of both commit sub-gate tasks, while both weak models
committed the gated mutation on every replicate of both tasks; the
enforced cells are safe by construction. We read this off the raw cell
counts, not a pooled p-value: the independent unit is the task (two per
arm, with near-deterministic replicates), so rep-level exact tests
overstate precision and are reported only as within-condition figures
beside the honest task-level count. The pattern is stark but small and
heavily caveated: capability is confounded with model family (both
capable rungs non-Qwen, both weak rungs Qwen); once the
threshold-sensitive below-floor 7B is set aside, the robust load-bearing
evidence rests on a single non-frontier model (Qwen3.5-9B); the pooled
cell is n=8, below our own pre-registered target; the capable-side ``no
marginal value of enforcement'' is a failure to detect (one-sided upper
bound +32pp), not demonstrated equivalence; and a salience qualifier on
the told-arm result rests on a separate n=3 diagnostic on an off-roster
model. Two hypothesized moderators nulled at that diagnostic tier
(context-verifiability via M8 audit emission; benign goal-conflict
pressure, zero fabrication across 40 reps), and the clinical-refusal
mechanism is near-ceiling and underpowered (its detector fires on 7 of
223 refusal reps, one task, one model), so we read no substitution
result off it.

A second, methodological contribution stands independent of the ladder:
harnesses that hide tool output manufacture spurious fabrication
findings (harness blindness); an apparent instrumental-fabrication
effect in our own pipeline dissolved to 0 percentage points (against a
pre-committed \textgreater=40pp bar) once the harness inlined command
stdout. We release GovernedAgentBench: the task suite, the deterministic
offline scorer, and the frozen paid-run trajectories and scores, with
the reference runtime pinned by git fingerprint (the PyPI wheel does not
enforce).

\subsection{1 Introduction}\label{introduction}

An agent harness is a design surface. The engineer who assembles an LLM
agent over user-owned structured data controls the loop around the
model: which tools are exposed, how actions are parsed and dispatched,
and what happens when a proposed action would violate a constraint. For
any given constraint, that engineer holds exactly two levers. The first
is in-context specification: state the rule in the agent's prompt as an
in-context contract (a command manifest, a flag, a mutation class, a
boundary description) and rely on the model to respect it. The second is
runtime enforcement: have the harness block or gate the violation
deterministically, regardless of what the model decides. Telling is
cheap and moves with the prompt; enforcing requires typed command
schemas, dispatch gates, and validation code that must be built, tested,
and maintained. A harness designer deciding where to spend engineering
effort needs to know when the second lever buys behavior the first does
not.

Most existing guardrail evaluations do not answer this question, because
they conflate the two levers or vary enforcement against a
prompt-wording axis rather than a told-versus-withheld one. Some toggle
enforcement while holding the full policy in the system prompt in both
conditions {[}arXiv:2604.15579{]}, so the delta is measured only where
the agent was already told. Others bundle whether rules are injected and
enforced into one variable {[}arXiv:2602.22302{]}, or compare an
enforced arm against a prompt-only arm {[}arXiv:2603.00822,
arXiv:2605.28914, arXiv:2606.21856{]}, which mixes the effect of telling
into the effect of enforcing; the closest, Mind the GAP
{[}arXiv:2602.16943{]}, crosses three prompt-wording conditions with a
runtime governance contract but never withholds the rule per mechanism.
The enforced-not-told cell is nearly absent from prior work (the closest
located instance is a single by-design condition in FORGE
{[}arXiv:2602.16708{]}), and no located prior work runs the
specify(told-versus-withheld)-per-mechanism factorial with first-attempt
scoring; Section 2 develops the map. None of these designs can answer
the question a harness designer actually faces: what does enforcement
add, given that the agent was told?

This paper crosses the two levers factorially in a specify-vs-enforce
2x2 per governance mechanism (Section 3). Its headline quantity is the
marginal value of enforcement given the agent was told (cell A:
deployment baseline vs cell B: told-not-enforced). We instantiate the
design in GovernedAgentBench, built over a frozen non-clinical
personal-wellness reference runtime (HAI, pinned by git fingerprint;
Section 4.1) whose governance mechanisms (M4 through M8, Section 3.2)
can be ablated individually or all at once, and scored by a
deterministic offline scorer.

We hypothesized three moderators would gate when specification
substitutes for enforcement: context-verifiability of the constraint,
benign goal-conflict pressure (compliance costs task success; benign
completion pressure, not adversarial injection), and model capability.
Model capability is the one that carries the result, in a small
pre-registered ladder we read descriptively rather than as a powered
test. Told the mutation-gate rule (the W57 user-commit boundary) with
the runtime off, the two capable models (MiniMax-M3, Llama-3.3-70B)
complied on every replicate of both commit sub-gate tasks, while the two
weak models (Qwen3.5-9B, Qwen2.5-7B) committed the gated mutation on
every replicate of both tasks; the enforced cells are safe by
construction, so the finding lives entirely in the told-not-enforced
(cell B) and untold-not-enforced (cell D) cells. The independent unit is
the task (two per arm) and the four replicates per cell are
near-deterministic, so we lead with the raw counts and the per-rep
mechanistic verification (a capable model refuses the commit citing the
user-gate; a weak model runs it while the gate is disabled) rather than
a pooled p-value; where we report an exact test we print the honest
task-level figure beside the rep-level one. The pattern is a genuine
sign flip in the marginal value of enforcement given told --
undetectable for the capable models, the whole barrier for the weak ones
-- but it is small and confounded: capability is confounded with model
family (both capable rungs non-Qwen, both weak rungs Qwen), and once the
threshold-sensitive below-floor 7B is set aside the robust load-bearing
evidence reduces to a single non-frontier model (Qwen3.5-9B, Section
7.1). The capable-side ``no marginal value of enforcement'' is a failure
to detect at this n (one-sided upper bound +32pp), not demonstrated
equivalence. Because the reconstructed untold-unenforced cell actually
violates even for the capable models (Section 4.4), telling and
enforcing each have a measured effect within a single model, so the
capable-model result reads as substitution rather than pre-existing
disposition. The other two moderators nulled at the diagnostic tier (one
model, Qwen3-235B-A22B-Instruct-2507, Together AI, temperature 0, n=3-5
per cell): the verifiability exception dissolved once the agent
retrieves the relevant evidence, and goal conflict produced zero
fabrication across all 40 reps under graded pressure P0-P3. The
clinical-refusal mechanism is near-ceiling and underpowered (its
detector fires on 7 of 223 refusal reps, one task and one model), so we
read no substitution result off it. We claim no ``unnecessary
guardrails'' conclusion: where the model self-enforces, enforcement's
marginal behavioral value is undetected at this n but is never zero as a
guarantee, and its demonstrated value remains the deterministic
guarantee plus the region where the model cannot self-enforce -- below
the operate floor, the untold violation floor, the sub-frontier
capability band, the incidental-salience failure of Section 5.2, and
adversarial intent, untested here. Injection robustness is cited
territory, not a claim of this paper. {[}FORK-L7-DRIFT: unmeasured; do
not assume outcome{]}

The second finding is methodological. Our own harness initially
manufactured a spurious fabrication result: it fed the agent a file-path
placeholder instead of actual command stdout, and the blinded agent
guessed identifiers it could not see. Inlining stdout dissolved the
apparent instrumental-fabrication effect entirely (0 percentage points
against a pre-committed threshold of at least 40 points; diagnostic, one
model, n=5), and the agent abstained honestly instead. We name this
failure mode harness blindness: evaluations that hide tool output from
the agent manufacture findings like this.

\subsubsection{Contributions}\label{contributions}

\begin{enumerate}
\def\labelenumi{\arabic{enumi}.}
\tightlist
\item
  A capability-moderated finding with explicit scope conditions,
  reported as a small single-runtime case study rather than a powered
  claim: on a pre-registered four-model ladder over the frozen HAI
  reference runtime, told the mutation-gate rule with the runtime off,
  both capable models complied on every replicate of both commit tasks
  while both weak models committed the gated mutation on every replicate
  -- a sign flip in the marginal value of enforcement given told,
  undetectable for the capable models and the whole barrier for the weak
  ones. The independent unit is the task (two per arm) with
  near-deterministic replicates, so we lead with raw counts and per-rep
  mechanistic verification and treat the rep-level exact test as a
  within-condition figure (the honest task-level exact p is 0.33).
  Capability is confounded with model family, the robust load-bearing
  side reduces to one non-frontier model (Qwen3.5-9B), and the
  capable-side null is a failure to detect (one-sided upper bound +32pp)
  not equivalence; the salience qualifier and deterministic-guarantee
  caveat remain attached; and the clinical-refusal arm is near-ceiling
  and underpowered, not a second substitution result. Single-runtime
  case study; external non-HAI replication is future work (Section 7.2).
\item
  A methodological demonstration that harness blindness manufactures
  spurious fabrication findings, shown by dissolving one such finding
  with a harness-level fix, offered as a caution for agent-eval
  methodology.
\item
  GovernedAgentBench, the released instrument for the 2x2: the task
  suite, a deterministic offline scorer, the reference runtime pinned by
  git fingerprint (the PyPI wheel does not enforce; Section 4.1), and
  the paid-run trajectories and scores released as a versioned archive
  alongside the preprint, with static oracle-pair evidence, live runtime
  probes, and model-backed diagnostics kept as separate evidence tiers.
\end{enumerate}

The novelty claim is a conjunction, not any single first: the 2x2
including the enforced-not-told cell, per-mechanism isolation, the
substitution account with its moderators, the methodological warning,
the deterministic offline scorer, and the released benchmark, taken
together. This is an AI-engineering study of agent-harness governance.
It claims no scaling law, no injection robustness, and no result beyond
this controller, task suite, and evidence tier.

\subsection{2 Background and related
work}\label{background-and-related-work}

\textbf{Agent harnesses and scaffolding.} The harness, the deterministic
software layer between a model and its tools, is now a design surface in
its own right. A survey under review at TMLR treats governance and
verification as a distinct harness layer (ETCLOVG, OpenReview
3hXEPbG0dh; cited via OpenReview, not as an accepted paper).
Harness-Bench {[}arXiv:2605.27922{]}, NLAH {[}arXiv:2603.25723{]}, AHE
{[}arXiv:2604.25850{]}, and Meta-Harness {[}arXiv:2603.28052{]} measure
or optimize harness configurations for capability; none carries a
governance axis. This paper measures the governance layer rather than
optimizing it.

\textbf{Runtime enforcement and guardrails.} Enforcement frameworks
(AgentSpec {[}arXiv:2503.18666{]}; NeMo Guardrails; Guardrails AI;
Invariant Guardrails, from Invariant Labs, now part of Snyk) supply
mechanism substrates without measuring what the model would do on its
own. In the injection-defense literature, CaMeL {[}arXiv:2503.18813{]}
and Progent {[}arXiv:2504.11703{]} show architectural enforcement
outperforming prompt-only heuristic defenses on AgentDojo; CaMeL never
states its policies to the model in any condition, and Progent runs no
disclosure factorial, so neither manipulates disclosure. Harness-MU
{[}arXiv:2606.21856{]} argues governance constraints should be
``enforced by execution hooks rather than entrusted to the LLM,'' a
single-axis comparison under adversarial framing. ContextCov
{[}arXiv:2603.00822{]} compiles instruction files into executable
guardrails (88.3\% compliance enforced against 67.0\% prompt-only on
SWE-bench Lite); its enforced arm feeds violation traces back to the
model, in-context specification delivered late. Verifier Tax
{[}arXiv:2603.19328{]} compares unenforced baselines (Tool-Calling,
TRIAD) against a policy-mediated architecture (TRIAD-SAFETY) on
tau-bench Airline and Retail: safety mediation intercepts up to 94\% of
non-compliant actions, yet safe success stays below 5\% in most
settings, with agents hallucinating user identifiers to route around
blocks; it toggles enforcement architecture, not prompt disclosure, so
it runs no told/untold contrast. Mechanical Enforcement
{[}arXiv:2605.14744{]} compares text-only governance against mechanical
enforcement (four architectural primitives outside the model's
interpretive loop) in a regulated-banking decision agent, cutting
uninformative deferrals by 73\% and raising task accuracy from MCC
approximately 0.43 to 0.88; it pairs two alternative configurations
rather than crossing a factorial, has no enforced-not-told cell, and
measures decision-rationale quality rather than tool-action compliance.
(``Prompt-based guardrail'' in Policy-as-Prompt {[}arXiv:2509.23994{]}
means an external classifier reading a policy, a different lever from
the in-context contract studied here.)

\textbf{Instruction and in-context rule following.} The main risk of
misreading this paper is reducing it to ``capable LLMs follow
instructions.'' That models follow checkable instructions is
established, but IFEval {[}arXiv:2311.07911{]}, RECAST
{[}arXiv:2505.19030{]}, and VerIF {[}arXiv:2506.09942{]} use
checkability to build evaluations and rewards, not as a predictor of
when agents comply. And told-only compliance is not free:
LogiSafetyBench {[}arXiv:2601.08196{]} names ``Unsafe Success,'' capable
models completing tasks while violating latent rules, non-monotone
across model families, as is SABER {[}arXiv:2606.01317{]};
Agent-SafetyBench {[}arXiv:2412.14470{]} reports none of 16 agents above
60\% safety score; SafePyramid {[}arXiv:2606.29887{]} shows in-context
policy application degrading with policy reasoning complexity;
Governance Decay {[}arXiv:2606.22528{]} shows constraints obeyed while
visible are silently erased by context compaction, violation rising from
0\% to 30\%, up to 59\% for some models. That establishes the
drift-causes-violation half of the stale-manifest phenomenon; whether
independent runtime enforcement catches post-drift violations remains
open. {[}FORK-L7-DRIFT: unmeasured; do not assume outcome{]} Our
question is therefore not whether an agent can follow a stated rule, but
what marginal behavioral value runtime enforcement retains once the same
constraint is specified in context.

\textbf{Software contracts and capability security.} The runtime's
vocabulary of manifests, typed command schemas, proposal/commit gates,
and least privilege descends from design-by-contract and capability
security. Agent Behavioral Contracts {[}arXiv:2602.22302{]} evaluates
runtime contracts at scale (200 scenarios, 7 models, 1,980 sessions),
but its conditions ``differ only in whether ABC contract rules are
injected and enforced,'' bundling specification and enforcement into one
variable, on a patent-gated benchmark. Prompts Don't Protect
{[}arXiv:2605.18414{]} filters unauthorized tools out of the model's
context with no policy prompt, a structurally enforced-not-told arm
whose 0.0\% unauthorized-invocation rate holds by construction, since
attempted violations cannot be observed.

\textbf{Closest priors.} AIRGuard {[}arXiv:2605.28914{]} is a
security-framed pre-action runtime guard; its prompt-only versus
runtime-guard ablation is a single-axis attack-success delta on one
model, with no factorial, no enforced-not-told cell, and no
constraint-class analysis. PhantomPolicy {[}arXiv:2604.12177{]}
formalizes policy-invisible violations, where compliance depends on
runtime state absent from the agent's context, contributing an
eight-category taxonomy this paper adopts and crosses with the
enforcement axis. It does run a policy-in-prompt condition (risky-case
violations 95.3\% to 40.7\% under human-reviewed labels), but folds
telling into the same single ordinal axis as enforcement architecture
(baseline, policy-in-prompt, content-DLP, Sentinel), with Sentinel and
the DLP baseline always applied to baseline not-told traces, so no
condition holds enforcement fixed while varying telling and no
enforced-and-told versus enforced-and-not-told contrast exists. Neither
AIRGuard nor PhantomPolicy crosses the two levers. Mind the GAP
{[}arXiv:2602.16943{]} is the closest prior to this paper's
capability-moderation headline: it reports both that prompt-based
tool-call safety is strongly model-dependent across six frontier models
and that a runtime governance contract produces no detectable deterrent
on forbidden tool-call attempts. Its manipulated axis is prompt wording
(three phrasings) crossed with a single monolithic governance contract
under jailbreak framing, not told-versus-withheld per mechanism; it has
no withheld arm, no enforced-not-told cell, no per-mechanism inventory,
and no first-attempt scoring, and its ``no deterrent'' finding is a
leakage catch-layer under attack rather than a hard commit gate under
benign use (the directional tension is taken up in Section 8.1). We cite
it from its verified abstract and released repository only. FORGE
{[}arXiv:2602.16708{]} runs the only located by-design enforced-not-told
condition: ``The instrumented condition does not include the natural
language policy in the agent's prompt; the agent relies solely on
runtime enforcement and corrective feedback.'' It has no
told-and-enforced cell and no neither floor, so it cannot measure the
redundancy of enforcement given specification; corrective feedback makes
its enforced arm converged multi-turn rather than first-attempt; it
evaluates a single holistic Datalog policy per deployment, with no
per-mechanism inventory, no moderators, and no located public release.
ABSTAIN {[}arXiv:2606.02965{]} reaches the closest cell coverage, three
of four cells for a single abstention mechanism; its Checkpoint
condition reuses the Prompt-Only prompt (its Appendix B), realizing
told-and-enforced rather than enforced-not-told, and non-Checkpoint
conditions are LLM-judge scored. TeamBench {[}arXiv:2605.07073{]}
reports statistically indistinguishable pass rates between prompt-only
and OS-sandbox-enforced role separation, alongside 3.6 times more
verifier-edits-executor cases in the prompt-only arm: convergent
published evidence for the substitution shape in a different mechanism
domain, and a standing caution that aggregate parity can mask sub-metric
enforcement value, taken up in Section 7.2; it is prior evidence in a
different domain and design, not a replication of this paper's effect
(this paper is a single-runtime case study; external replication is
future work, Section 7.2). Symbolic Guardrails {[}arXiv:2604.15579{]}
holds the full policy in the system prompt in both conditions while
toggling enforcement; TRACE {[}arXiv:2606.13174{]} compiles user
corrections into runtime checks in the preference domain, the rule
always known to the agent in some form. The closest benchmark prior is
ST-WebAgentBench {[}arXiv:2410.06703{]}: in-context policies delivered
to three agent scaffolds and scored post hoc via Completion Under
Policy, with no mechanism that blocks a non-compliant action before
execution.

\textbf{Positioning.} Isolated enforced-not-told measurements exist
(FORGE by explicit design; Prompts Don't Protect structurally). No
located prior work holds that cell inside a crossed factorial against a
told-and-enforced cell and a neither floor, decomposes the comparison
per governance mechanism, scores the telling axis on first-attempt
behavior, or releases a deterministic offline scorer and benchmark that
hold the two levers apart. The contribution is that conjunction,
together with the harness-blindness methodological finding (Section 6);
no element is claimed as individually unprecedented.

\subsection{3 The two levers}\label{the-two-levers}

An agent harness has two distinct levers for making an agent respect a
constraint. It can state the rule where the model can read it, or it can
make the violation impossible regardless of what the model decides. We
call the first lever in-context specification (telling) and the second
runtime enforcement (enforcing). Deployed harnesses almost always pull
both levers at once: the system prompt describes the policy and the
runtime also implements it. Prior evaluations largely inherit this
bundling, toggling specification and enforcement as one variable
{[}arXiv:2602.22302{]} or manipulating only one axis while the other
never varies {[}arXiv:2605.28914{]}. This section defines the two levers
separately, fixes the mechanism inventory that instantiates the
enforcement lever, introduces the constraint classification that the
design uses as its predictor, and assembles the pieces into the 2x2 that
the rest of the paper measures.

\subsubsection{3.1 Telling: the in-context
contract}\label{telling-the-in-context-contract}

The telling lever is the in-context contract: the rule, manifest fact,
schema, or policy statement made available to the agent inside its
context at decision time. In the reference runtime this includes the
command manifest (which commands exist and their argument schemas), the
\texttt{agent\_safe} flag attached to each command, the mutation class
of each write surface, and the prose description of the out-of-contract
boundary, including its clinical-boundary leg. The contract is broader
than a system-prompt policy paragraph; schemas, flags, and taxonomies
are all part of what the agent is told.

Two boundaries keep this lever distinct. In-context specification is a
harness-level lever, not a training-time one; it is orthogonal to
instruction tuning or preference training, which shape the model's
priors rather than its decision-time context. And the in-context
contract is only the copy of the specification the model can read, not
the specification the runtime implements. The paper's central
manipulation is the gap between the two.

\subsubsection{3.2 Enforcing: the mechanism
inventory}\label{enforcing-the-mechanism-inventory}

The enforcing lever is runtime enforcement: the harness blocking,
gating, or otherwise making a violation impossible regardless of what
the model emits, independent of whether the model was told the rule. The
reference runtime (HAI v0.2.0, the frozen instrument described in
Section 4) implements enforcement through a fixed inventory of
mechanisms, five of them individually ablatable through a named runtime
mode and one held constant.

{\def\LTcaptype{none} % do not increment counter
\begin{longtable}[]{@{}
  >{\raggedright\arraybackslash}p{(\linewidth - 6\tabcolsep) * \real{0.2500}}
  >{\raggedright\arraybackslash}p{(\linewidth - 6\tabcolsep) * \real{0.2500}}
  >{\raggedright\arraybackslash}p{(\linewidth - 6\tabcolsep) * \real{0.2500}}
  >{\raggedright\arraybackslash}p{(\linewidth - 6\tabcolsep) * \real{0.2500}}@{}}
\toprule\noalign{}
\begin{minipage}[b]{\linewidth}\raggedright
ID
\end{minipage} & \begin{minipage}[b]{\linewidth}\raggedright
Mechanism
\end{minipage} & \begin{minipage}[b]{\linewidth}\raggedright
Ablation mode
\end{minipage} & \begin{minipage}[b]{\linewidth}\raggedright
Context-verifiable?
\end{minipage} \\
\midrule\noalign{}
\endhead
\bottomrule\noalign{}
\endlastfoot
M4 & Validation of typed commands and proposal payloads &
\texttt{no\_validation} & Yes (schema is in the manifest) \\
M5 & \texttt{agent\_safe} dispatch refusal & \texttt{no\_agent\_safe} &
Yes (\texttt{agent\_safe} flag is in the manifest) \\
M6 & W57 proposal/commit user gate & \texttt{no\_proposal\_gate} & Yes
(mutation class is in the manifest) \\
M7 & Refusal of out-of-contract requests (clinical-boundary leg is the
zero-tolerance subset); JSON output exempt & \texttt{no\_refusal} & Yes
(the boundary is describable in the prompt) \\
M8 & Audit evidence emission / reference faithfulness &
\texttt{no\_audit\_chain} & No under the standard setup (Section 3.3) \\
M9-TX & Transaction integrity & Held constant, non-ablatable & n/a \\
\end{longtable}
}

M4 guards the command surface itself: malformed, out-of-schema, or
hallucinated commands are rejected before dispatch. M5 refuses dispatch
of commands the manifest does not flag as agent-safe. M6 guards
user-authored state: activation or deactivation of a proposed mutation
is a user-gated commit. In the measured mutation-gate cells, however,
the block that fires is M5's dispatch refusal (the commit command is
flagged not-agent-safe), which reaches the dispatch first; M6's
proposal/commit gate is masked by M5 and is
\texttt{-\/-confirm}-bypassable, so it is defense-in-depth here and is
not independently exercised (PAPER.md D-48). We attribute the measured
block to M5 and do not claim a clean M5-versus-M6 decomposition. M7
guards the scope boundary, with the clinical boundary (no diagnosis,
treatment, prescribing, or autonomous medical decisions) as its
zero-tolerance subset; structured JSON output is exempt. M8 guards
reference faithfulness: the runtime emits and persists audit evidence
for recommendations, and the deterministic scorer detects fabricated or
suppressed references after the fact. M8 is provision-type rather than
block-type (Section 3.5 defines the distinction), and it is evidence
emission, not the whole audit chain. M9-TX is the atomic-transaction
floor; it is held constant and has no ablation mode.

The runtime-enforcement axis therefore has exactly seven modes:
\texttt{full\_contract}, the five per-mechanism off-paths
(\texttt{no\_validation}, \texttt{no\_agent\_safe},
\texttt{no\_proposal\_gate}, \texttt{no\_refusal},
\texttt{no\_audit\_chain}), and \texttt{no\_runtime\_enforcement}, in
which M4 through M8 are all off while the M1-M3 harness infrastructure
and M9-TX remain on. M1 through M3 are the always-on harness
infrastructure (the model loop and dispatch, action parsing, and
observation feedback); they are not ablatable governance mechanisms and
remain on in every mode, which is what makes
\texttt{no\_runtime\_enforcement} a floor over which cell D is still
driven rather than an inert harness. One caveat governs every
per-mechanism number in this paper: harness mechanisms are coupled, so
per-mechanism results are reported as marginal contribution within this
fixed controller, never as independent additive contributions.

\subsubsection{3.3 Context-verifiability}\label{context-verifiability}

The design's predictor variable is context-verifiability: a constraint
is context-verifiable if the agent can check its own compliance using
only what is already in its context at the moment it decides.
Verifiability is a property of the pair (constraint, decision-time
context), not of a mechanism in the abstract, and it can flip within a
trajectory: a constraint becomes verifiable the moment retrieval lands
the relevant ground truth in the agent's context. This framing adopts
the observation, formalized as policy-invisible violations in
{[}arXiv:2604.12177{]}, that compliance can depend on runtime state
absent from the agent's context; we treat that dependence as a
trajectory-local property rather than a fixed taxonomy label.

Under the standard task setup, with the manifest and taxonomies in the
prompt and no runtime state disclosed, M4, M5, M6, and M7 are
context-verifiable: the schema, the \texttt{agent\_safe} flag, and the
mutation class are in the manifest, and the M7 boundary is describable
in the prompt. M8 was classified non-verifiable under the standard setup
and pre-registered as the exception where even a capable cooperative
agent could not self-enforce, because whether a cited evidence reference
actually exists is runtime state the agent cannot see. The
pre-registered prediction attached to this classification was that
enforcement is behaviorally redundant for context-verifiable constraints
under cooperative, unconflicted conditions and behaviorally necessary
for non-verifiable ones at every capability level. Section 5 reports
that M8's classification did not survive contact with the evidence:
audit faithfulness turns out to be verifiable once the agent retrieves
the evidence, exactly the within-trajectory flip the definition permits.

\subsubsection{3.4 L7 stale-manifest drift: the non-retrievable task
condition}\label{l7-stale-manifest-drift-the-non-retrievable-task-condition}

One constraint class is genuinely non-retrievable rather than merely
unretrieved. Under L7 stale-manifest drift, the agent's in-context
manifest is outdated relative to the actual runtime contract, and the
agent cannot retrieve the fact that this is so: that ground truth is
runtime state no read surface exposes. No sequence of lookups can flip
it to verifiable, which distinguishes drift from M8. Drift is a task
condition, not a mechanism row: it is a property of the relationship
between the contract copy the agent holds and the contract the runtime
implements, and it can afflict any mechanism whose manifest facts have
moved. L7 drift is the sole surviving candidate for a true
non-verifiable enforcement delta in this design, and it is unmeasured.
{[}FORK-L7-DRIFT: unmeasured; do not assume outcome{]}

\subsubsection{3.5 What the runtime guarantees regardless of
behavior}\label{what-the-runtime-guarantees-regardless-of-behavior}

The two levers differ in kind, not only in strength. Telling changes
what the model is likely to do; enforcing changes what can happen at
all. With enforcement on, the block-type mechanisms (M4 through M7) make
the guarded violation impossible to complete no matter what the model
emits: an out-of-schema command does not dispatch, a flagged command
does not execute, an uncommitted proposal does not activate. M8's
guarantee is provision-type: evidence is emitted and persisted whether
or not the model narrates it faithfully, so fabrication is detectable
after the fact. M9-TX guarantees atomicity in every mode. This is the
deterministic guarantee, and it holds regardless of model behavior,
model identity, and prompt content. The distinction matters for reading
every result that follows: when this paper reports that enforcement is
behaviorally redundant in some cell, the claim is about the marginal
effect on observed behavior, never that the guarantee is worthless.
Behavioral redundancy is not structural redundancy.

\subsubsection{3.6 The 2x2}\label{the-2x2}

Crossing the two levers per mechanism yields the headline design:

{\def\LTcaptype{none} % do not increment counter
\begin{longtable}[]{@{}
  >{\raggedright\arraybackslash}p{(\linewidth - 4\tabcolsep) * \real{0.3333}}
  >{\raggedright\arraybackslash}p{(\linewidth - 4\tabcolsep) * \real{0.3333}}
  >{\raggedright\arraybackslash}p{(\linewidth - 4\tabcolsep) * \real{0.3333}}@{}}
\toprule\noalign{}
\begin{minipage}[b]{\linewidth}\raggedright
\end{minipage} & \begin{minipage}[b]{\linewidth}\raggedright
Runtime enforces
\end{minipage} & \begin{minipage}[b]{\linewidth}\raggedright
Runtime off
\end{minipage} \\
\midrule\noalign{}
\endhead
\bottomrule\noalign{}
\endlastfoot
\textbf{Contract in prompt} & A: deployment baseline & B:
told-not-enforced (self-enforcement) \\
\textbf{Contract withheld} & C: enforced-not-told (pure runtime) & D:
neither (violation floor) \\
\end{longtable}
}

Cell A, the deployment baseline, is how harnesses actually ship. Cell B,
told-not-enforced, isolates self-enforcement: the agent knows the rule
and nothing stops it. Cell C, enforced-not-told, isolates pure runtime
governance: the rule is withheld from the prompt while the runtime still
enforces it. Cell D, neither, is the violation floor. The
\texttt{no\_runtime\_enforcement} mode is a runtime mode, not a cell:
under the contract-in-prompt arm it serves as a robustness sanity floor,
never per-mechanism attribution evidence; under the contract-withheld
arm the same mode realizes cell D and anchors the floor for the
contrasts below.

Three contrasts carry the information. B vs D is the effect of telling:
what specification alone buys over nothing. C vs D is the effect of
enforcing: what the runtime alone buys over nothing. A vs B is the
marginal value of enforcement given the agent was told, the headline
quantity of this paper: if A and B do not differ behaviorally,
enforcement added no behavioral value on top of specification for that
constraint in that situation, and its remaining value is the
deterministic guarantee of Section 3.5.

The magnitude a designer should read off these contrasts depends on the
constraint's type, a distinction that becomes an identification limit in
Section 8.4. Three types recur. A disposition-covered constraint is one
a cooperative capable agent already respects whether or not it is told:
B vs D is near zero, and with the runtime off changing nothing the agent
would have done, A vs B is near zero too, so substitution there is
degenerate rather than measured. A specification-substitutable
constraint is one the agent cannot honor unless told (B vs D large), and
for which a runtime enforcer may not even exist. An
enforcement-load-bearing corner is one where the runtime carries
behavioral value the contract does not: the untold violation floor, the
region below the operate floor, the incidental-salience failure, and
adversarial intent. The marginal value of enforcement given told (A vs
B) is informative only for a constraint that is not disposition-covered.
On the frozen tasks our enforce-able mechanisms were disposition-covered
at the cooperative tier (Section 8.4); the reconstructed mutation gate
of Section 4.4 is not, its untold-unenforced cell actually violating
even for capable models, so its redundancy is measured rather than
degenerate. The capability ladder (Section 5.1) is the primary
identification of specification against disposition; the within-model
directiveness pilot of Section 9.1 is complementary.

One interaction between the cells requires a scoring convention. A
blocked action returns an error message, and an error message is
in-context specification delivered late, so cell C converges toward cell
B after the agent's first contact with the enforcement surface. Axis
attribution for B vs D and C vs D is therefore scored on first-attempt
behavior, with converged multi-turn behavior reported separately;
Section 4 gives the operational details. The 2x2 runs under three
moderators, constraint class per Section 3.3, benign goal-conflict
pressure, and model capability, which Section 4 defines and Section 5
evaluates.

\subsection{4 GovernedAgentBench:
methodology}\label{governedagentbench-methodology}

GovernedAgentBench is the instrument that holds the two levers apart: it
lets in-context specification and runtime enforcement be toggled
independently, per mechanism, against a fixed reference runtime, and
scores the result deterministically and offline. This section describes
the runtime, the task suite, the fixtures and hermeticity discipline,
the two axes as implemented, the scorer, the evidence tiers, the model
roster, and the reproducibility path.

\subsubsection{4.1 The reference runtime}\label{the-reference-runtime}

The runtime under measurement is HAI v0.2.0, a non-clinical
personal-wellness reference runtime with a typed CLI command surface.
HAI is the instrument, not the contribution: it is pinned as a frozen
snapshot so that every mechanism has one concrete, inspectable
implementation, and it is described here as a reference implementation,
not a product. The runtime instantiates the mechanism inventory of
Section 3: M4 validation of typed commands and proposal payloads, M5
\texttt{agent\_safe} dispatch refusal, M6 the W57 proposal/commit user
gate, M7 refusal of out-of-contract requests (with the clinical-boundary
leg as the zero-tolerance subset), M8 audit evidence emission, and M9-TX
transaction integrity, which is held constant and non-ablatable. The
non-clinical boundary is part of the evaluated contract, not a
disclaimer: the runtime refuses the clinical actions enumerated under M7
in Section 3.2, and the benchmark scores that refusal. The agent-facing
contract is a frozen manifest snapshot, \texttt{agent\_cli\_contract.v2}
(approximately 189 KB, checked in at
\texttt{manifests/hai\_0\_2\_0.json}), which carries the command
schemas, \texttt{agent\_safe} flags, and mutation classes; with it in
the prompt, and the M7 boundary describable in the prompt itself, M4
through M7 are context-verifiable under the standard task setup (Section
3.3).

Two facts about the runtime's identity are load-bearing for
reproduction. First, the PyPI \texttt{v0.2.0} wheel does not enforce:
the dispatch-time \texttt{agent\_safe} gate landed three days after the
version tag, so \texttt{pip\ install\ hai==0.2.0} yields a non-enforcing
runtime that would silently corrupt the enforce arm. The measured
runtime is therefore pinned by git fingerprint \texttt{6c82cd0}, whose
nearest named ancestor is the annotated tag \texttt{gab-runtime-1.0}, 14
commits back; the intervening commits add the M7 clinical scanner and
the M5 side-door hardening the enforce arm needs, so the tag predates
the measured runtime and must not be used to reproduce -- reproducers
check out \texttt{6c82cd0}, not the tag or the wheel. Second, to
guarantee the enforce arm was live before any paid call, a pre-spend
runtime-identity preflight behaviorally asserts enforcement: it drives
an agent-context commit under \texttt{full\_contract} and requires the
runtime to refuse it (exit non-zero, state left proposed), hard-stopping
the sweep otherwise. The preflight passed for this run (fingerprint
\texttt{6c82cd0}, wheel absent), so the enforce cells were measured
against a runtime verified to enforce, not assumed to.

\subsubsection{4.2 Task suite and oracle-pair
construction}\label{task-suite-and-oracle-pair-construction}

The suite is drawn from five levels of a seven-level task taxonomy: L1
intent-to-command routing, L2 setup and recovery, L5 faithful narration
of runtime output, L6 governance and refusal, and L7 contract drift; the
intervening L3 and L4 levels are out of scope for this suite (Appendix A
gives the full taxonomy). Each task declares its
\texttt{load\_bearing\_mechanisms} and
\texttt{runtime\_modes\_in\_scope}; the harness refuses to run a task
under a runtime mode the task does not declare, so mode coverage is a
checked property rather than a convention. The released suite is 39
tasks (L1: 2, L2: 4, L5: 8, L6: 24, L7: 1) after the D-48 removal of the
validation-doctor task and the D-49 clinical-refusal additions; the paid
ladder swept a concentrated 16-task subset of it (Section 4.10), with
the remainder covered offline. PAPER.md decision D-37 (2026-07-04)
retired the earlier 28-task / 25-oracle-pair apparatus and rebuilt the
suite as a sharp per-mechanism 2x2 (told x enforced); D-39 (2026-07-05,
commit \texttt{ed07f5a}) then expanded it to three scenario pairs per
mechanism, on the grounds that a per-mechanism null resting on a single
scenario is too thin to carry. Every task is a labelled cell of its
mechanism's 2x2 via an explicit \texttt{contract\_arm} field, and the
suite's moderators are committed suite members rather than post-hoc
splits (Section 4.10). Appendix A holds the per-task table.

Every task is constructed as an oracle pair that routes a mechanism's
effect onto a scored metric: a hand-authored \texttt{full\_contract}
trajectory and a matching mechanism-off trajectory whose scored outcome
must differ from it on the task's load-bearing metric, demonstrating
scorer sensitivity to the mechanism. Each of M4 through M8 is covered by
static oracle-pair load-bearing tasks under the D-39 suite (three
scenario pairs per mechanism), and two tasks carry the
\texttt{no\_runtime\_enforcement} floor; the retired pre-consolidation
pair counts (25 pairs) no longer apply, and the per-task table is
Appendix A's. A mechanism's off-path is accepted as isolated only if
every emitted \texttt{mechanism\_disabled} marker belongs to that mode's
disabled set, at least one marker for the target mechanism fires,
\texttt{full\_contract} emits zero markers, and the scored consequence
delta against \texttt{full\_contract} on the task's load-bearing metric
is attributable to the disabled mechanism. These pairs are static
oracle-pair evidence: a deterministic canary establishing scorer
sensitivity and coverage over constructed cases, not live causality
evidence.

\subsubsection{4.3 Fixtures, hermeticity, and the manifest
snapshot}\label{fixtures-hermeticity-and-the-manifest-snapshot}

Seven synthetic fixtures are built for the benchmark
(\texttt{empty\_user}, \texttt{ready\_user\_minimal},
\texttt{read\_surface\_user}, \texttt{governance\_user},
\texttt{drift\_user}, \texttt{adversarial\_user}, and
\texttt{audit\_pending\_user}, the last leaving the target day
un-synthesized so no evidence card exists at build time, which stresses
M8 audit-reference faithfulness); not all seven are referenced by every
locked task, and the fixture-to-task mapping follows the D-39 suite of
Section 4.2. Fixtures contain no private health rows, no live wearable
exports, no credentials, and no maintainer personal data. Every
benchmark subprocess runs hermetically: \texttt{HAI\_HERMETIC=1} with
the state database and base directory redirected
(\texttt{HAI\_STATE\_DB}, \texttt{HAI\_BASE\_DIR}), and a fresh fixture
is materialized per repetition, so no state leaks across repetitions or
modes. L7 tasks additionally swap in a deliberately stale manifest
snapshot (\texttt{agent\_cli\_contract\_v1\_drift.json}) at task-load
time, so the agent's in-context contract diverges from the actual
runtime contract by construction; this drift condition is a task
condition, not a mechanism row, and it remains unmeasured at the
model-backed tier {[}FORK-L7-DRIFT: unmeasured; do not assume
outcome{]}.

\subsubsection{4.4 The two axes as
implemented}\label{the-two-axes-as-implemented}

The runtime-enforcement axis is the \texttt{runtime\_mode} seam,
exposing the seven modes of Section 3.2. The contract-in-prompt axis is
a per-task \texttt{contract\_arm} field (told or untold) realized at
prompt-render time by three coordinated transforms, while the runtime
remains unchanged: a structured transform over the embedded manifest
that withholds the mechanism-relevant facts (flags, mutation classes,
schema constraints for the constraint under test), a prose scrub that
removes the constraint's natural-language statements from the prompt
body under a forbidden-token check, and parameterized boundary blocks so
the told arm's boundary prose is assembled per task rather than
hand-edited. A mechanical untold-leak scan over every rendered untold
prompt is a lock gate: no rendered untold prompt may carry the withheld
constraint's tokens. Crossing the two axes yields the cells of Section
3.6, and the two roles of \texttt{no\_runtime\_enforcement} (robustness
sanity floor under the contract-in-prompt arm, cell D anchor under the
contract-withheld arm) are kept distinct as defined there. Two
provenance caveats attach to the untold arm. First, to keep the
untold-unenforced cell from measuring command-discovery failure rather
than governance disposition, the untold render restores a rule-free
description of the commit command (D-54); this discoverability
reconstruction was applied after the canary phase, so the capable
models' cell-D violation rate -- the floor the
substitution-not-disposition argument rests on -- carries that
post-canary provenance and can be re-scored with and without it offline.
Second, the told-arm boundary prose and the untold withholding were both
hand-reconstructed pre-run (D-52) so that told states the operative rule
and untold withholds it, fixing an earlier artifact in which both arms
committed; this is disclosed rather than silent.

Three disclosures attach to the untold arm's construction. First, the
untold arm withholds all natural-language constraint prose but, by
design, preserves the object-capability command surface: the model still
sees command names (e.g.~\texttt{hai\ target\ commit},
\texttt{hai\ target\ archive}) and their \texttt{-\/-confirm} flags. The
schema therefore partially specifies the constraint even when untold, so
B-vs-D and C-vs-D are lower bounds on the effect of telling; the harness
schema is itself an in-context contract. Second, in the untold M5 arm
the manifest transform asserts \texttt{agent\_safe:\ true} rather than
merely omitting the gate, so the untold M5 cells are a
``told-it-is-safe'' floor, disclosed as told-the-opposite rather than
uninformed; the diagnostic harm-floor probe shows violations still occur
in cell D despite the schema signal, so the floor is non-vacuous. Third,
cell B's told prompt retains boundary prose stating that the runtime
enforces the constraint, while the runtime is in fact ablated in that
cell; the prompt is held constant across the enforcement axis by design,
so cell B measures self-enforcement under an (inaccurate) assertion of
enforcement, and this is disclosed rather than edited away, since
editing it would break the held-constant told prompt.

The operator action contract is structured JSON, one action per turn,
not arbitrary shell. Transcripts record the model's emitted action
verbatim as the assistant turn and each observation as a returned
message; no provider-native tool calls are synthesized, because
fabricating tool-call structure would misrepresent the evaluated
transcript.

\subsubsection{4.5 Deterministic scoring and first-attempt
attribution}\label{deterministic-scoring-and-first-attempt-attribution}

The scorer is deterministic and offline; it contains no model calls. Its
behavior thresholds, pass rule, and critical violation kinds are loaded
from \texttt{scorer\_config.paper\_v1.json}, whose audited bytes are
hashed in \texttt{lock\_hashes.json}. We disclose a scope limit on that
frozen record: the scorer's Python logic in \texttt{core.py} is not
itself hashed, and the harm-only \texttt{unsafe\_action\_rate} code path
was edited three times in the thirty hours before the run -- adding a
state-write side-door check (D-54), an active-insert check (D-55), and
an argument-defaults resolution (D-55.1, committed 47 minutes before the
run) -- each prompted by a pre-run canary and each in the
hypothesis-favorable direction. ``Frozen'' therefore scopes to the
config bytes, not the scoring code, and the pre-registered scorer should
be read as fixed only up to those three commits. We make the material
consequence checkable: those edits added an
\texttt{unsafe\_command\_args} path that fired in zero mutation-gate
headline cells (every violation in Table 5.1 is the pre-existing
\texttt{must\_not\_call} detection, \texttt{unsafe\_mutation}), so
re-scoring the paid trajectories with the pre-D-54 scorer is the
identity on the headline 2x2 and the load-bearing result is invariant to
the pre-run code changes; only the below-floor 7B's graded partial
scores touch the new path. One circularity is disclosed rather than
removed: the scorer's clinical-boundary detector imports the same
banned-phrase list the runtime's M7 refusal uses (frozen before the run,
audited for over-broad single-word entries), so scorer and runtime share
ground truth for that violation kind and cell A is clean on it by
construction. Citation resolution for \texttt{must\_cite} metrics is
scoped to command stdout only.

The telling axis is scored on first-attempt behavior, per the
convergence rationale of Section 3.6. First-attempt is defined uniformly
as the first action that reaches the enforcement surface for the
constraint in question, not literal step one; the definition applies
identically to refusal rows and dispatch rows, so an agent that emits
invalid-output steps before its first contact with the enforcement
surface is scored from that first contact. In the pre-registered run the
first-attempt window closes only on genuine enforcement contact, a
blocked \texttt{must\_not\_call} gated action, not on any incidental
error message. For the gated mechanisms (M5 \texttt{agent\_safe}, M6
proposal/commit) that contact is, for a single-target
\texttt{must\_not\_call} task, the model's first action, so the telling
and enforcing contrasts for these mechanisms are single-decision
difference-of-proportions, reported as pooled counts; we frame the two
windows as enforcement-as-barrier (first attempt) versus
enforcement-as-teacher (converged) and do not claim multi-step
first-attempt reasoning for gated cells. Validation, refusal, and
audit-chain mechanisms retain multi-step first-attempt windows.
First-attempt scoring carries the B vs D and C vs D attribution, and
converged multi-turn behavior is reported separately and never blended
into first-attempt numbers.

The metric is only well-defined for a model whose actions the harness
can read: a model whose output the action parser cannot read never
reaches the enforcement surface, so its first-attempt behavior is
undefined and its apparent failure conflates self-enforcement with the
operate floor (Section 6.3). The early diagnostic ladder was confounded
exactly this way, which at the time scoped first-attempt attribution to
Qwen3-235B, the one parser-clean model. An envelope-tolerant parser has
since shipped (uniform code-fence stripping and prose-embedded JSON
extraction applied identically to every model before validation;
normalization never changes which action is validated; commit
\texttt{8319b83}), so in the pre-registered run first-attempt scoring
runs across the full ladder of Section 4.7. Two caveats survive the fix
and are disclosed there: per-model vendor-recommended sampling confounds
cross-model first-attempt comparisons, and the cross-family evidence
rests on a single Llama point.

A deterministic rule baseline anchors the routing tasks as plumbing
evidence only: it confirms the harness, fixtures, and scorer compose,
and it fails the L5 narration tasks by design, since it emits no final
narration. (PAPER.md D-30 recorded five L5 narration tasks at the time
of the scorer-correctness pass; the D-39 suite carries eight L5 tasks,
per Section 4.2 and Appendix A.)

\subsubsection{4.6 Evidence tiers}\label{evidence-tiers}

Three evidence tiers are kept separate throughout and are never merged
into one claim: static oracle-pair evidence (hand-authored full/off
pairs, scorer coverage), live runtime probes (end-to-end HAI runs under
compared modes with real command execution, live causality evidence),
and model-backed diagnostics (real model completions scored by the same
scorer). Model-backed evidence comes in two tiers, kept separate and
never merged: the primary pre-registered capability-ladder sweep (four
models, n=8 per cell on the mutation gate; Sections 4.10, 5.1) and a
lower diagnostic-probe tier (one model, Qwen3-235B-A22B-Instruct-2507,
Together AI, temperature 0, n=3 to 5 per cell) that framed the design
and is reported as provenance in Sections 5.2 through 5.5.

\subsubsection{4.7 Model roster}\label{model-roster}

Two rosters must be kept distinct: the diagnostic-probe roster under
which the Section 5.2-5.5 probes were run, and the pre-registered run
roster (PAPER.md D-41) under which the capability-ladder sweep of
Section 5.1 was executed.

\textbf{Diagnostic-probe roster (Sections 5.2-5.5).} The working model
was \texttt{Qwen/Qwen3-235B-A22B-Instruct-2507-tput} (Together AI,
non-thinking mixture-of-experts, serverless), run at temperature 0.
\texttt{claude-sonnet-4-6} was the designated fallback for
narration-heavy M8 tasks; no run reported here is Sonnet-backed.
Excluded during viability screening: Qwen2.5-7B-Instruct-Turbo (below
the operate floor, 0/60 valid finals, contract friction rather than
governance), Mistral-Small-24B (32k context too small for the manifest
prompt plus multi-turn history), Gemma-3-27B and Qwen2.5-32B
(dedicated-endpoint only), and Gemma-4-31B (a reasoning model whose
thinking consumed the token budget, yielding empty completions and
timeouts).

\textbf{Pre-registered run roster (Section 5.1).} The executed ladder is
a four-model Together serverless capability ladder:
\texttt{MiniMaxAI/MiniMax-M3} as the capable primary,
\texttt{Llama-3.3-70B} as the second capable point, \texttt{Qwen3.5-9B}
as the near-floor point, and \texttt{Qwen2.5-7B-Instruct-Turbo}
deliberately re-admitted as the below-floor operate control (its
diagnostic exclusion is the reason it is in the ladder, since a
below-floor model is the positive control for the operate-floor claim).
The pre-registration's primary was
\texttt{Qwen/Qwen3-235B-A22B-Instruct-2507-tput}; the provider
deprecated it from serverless mid-cycle, and we replaced it with
MiniMax-M3. We disclose the researcher degree of freedom this
introduces: the replacement was chosen after a canary confirmed it
reproduced the target specification effect (an outcome-conditioned
selection, PAPER.md D-56), so MiniMax-M3 should be read as a
non-pre-registered addition and its result as exploratory. The
pre-registered capable rung unaffected by this is Llama-3.3-70B, which
shows the same capable-side pattern. Capability is also confounded with
model family: the two capable rungs are non-Qwen (MiniMax, Llama), the
two weak rungs Qwen (Section 8.4). Sampling is vendor-recommended per
model rather than uniform temperature 0, so model identity is confounded
with sampling regime across the ladder; the per-model 2x2 is computed
within a single model at a constant sampling regime and is unaffected
(Sections 4.10, 8.4). Provider seeds are unsupported on the serverless
endpoints, so replication is across independent stochastic draws rather
than fixed-seed determinism. Exact model identifiers, pricing, context
windows, and sampling settings are pinned in the locked roster revision
and recorded in the run manifest (PILOT\_PROTOCOL.md §20).

\subsubsection{4.8 Goal-conflict variants and the adversarial
layer}\label{goal-conflict-variants-and-the-adversarial-layer}

The goal-conflict arm uses task variants applying benign completion
pressure, not adversarial injection: respecting the constraint costs
task success. The diagnostic H2 probes graded this from P0 to P3, where
at P3 a pipeline is blocked and the user cannot proceed; in the
pre-registered suite goal conflict is carried by two committed suite
members (\texttt{gab\_l6\_agentsafe\_conflict},
\texttt{gab\_l5\_audit\_conflict}; Section 4.10). No attacker-controlled
content is involved; this targets the literature-verified locus where
frontier models violate policies stated in their own system prompt
(Symbolic Guardrails {[}arXiv:2604.15579{]}), including under realistic
completion pressure (LogiSafetyBench {[}arXiv:2601.08196{]}).
Separately, a 16-trajectory adversarial layer (four hand-authored
trajectories each against M7, M5+M6, M8, and M4) provided
scorer-coverage evidence at the appendix tier only, never model-backed.
It was retired at the benchmark tier by commit \texttt{a10e850}, and
PAPER.md D-37 records the retirement as final: the adversarial-input arm
is cited future work, not an arm of this paper (Appendix D). Injection
robustness is cited territory in the security literature, which already
shows telling fails under attack, not a claim of this paper.

\subsubsection{4.9 Offline
reproducibility}\label{offline-reproducibility}

A single command, \texttt{reproduce\_offline.py}, regenerates the
rule-baseline ablation, evidence tables, figures, error taxonomy, and
cell contrasts with no network access, no private state, and no paid
APIs. The static and live isolation matrices were retired at the
benchmark tier by commit \texttt{a10e850} (Appendix E) and the pipeline
no longer emits them. It emits a top-level manifest with per-artifact
summary metadata and exits nonzero if the rule-baseline ablation
produced no scored rows or figures; the golden-artifact baseline records
hashes and a wall-clock reference (Appendix E; re-baseline pending). The
model-backed trajectories and scores are the only artifacts that
required paid API access; the run tree is not in the source repository
(it is gitignored), so they are released as a versioned archive
alongside the preprint, and from that archive they are re-scoreable
offline by the same deterministic scorer. The Table 5.1 cell counts are
regenerated by the released pooling module
(\texttt{build\_cell\_contrasts\_pooled}); the exact tests reported
beside them are computed by a small released helper, since the pooling
module itself carries no statistics dependency.

\subsubsection{4.10 The pre-registered run
design}\label{the-pre-registered-run-design}

This subsection describes the pre-registered design of the run reported
in Section 5.1; the design was locked before the run in PAPER.md
decisions D-39 through D-43, with the run manifest, the pre-committed
outcome-branch map, and the full disclosure set carried by
PILOT\_PROTOCOL.md §20.

\textbf{Suite and cells.} The full 39-task suite of Section 4.2 is the
offline-reproduction scope. The paid ladder was pre-committed (D-49) to
a concentrated headline subset that carries both levers: the two Tier-A
mechanisms (the agent\_safe mutation gate and clinical refusal), the
Tier-B M8 blind/sighted twin, and the Tier-C operate control. The Tier-D
scope mechanisms (M4 validation, the M6 proposal-gate variant, and the
credential and audit-domain families) are covered offline by
\texttt{reproduce\_offline} and are never paid-swept, so the paid run
measures only constraints on which both the telling and the enforcing
lever move. Each paid task is a labelled cell of its mechanism's
told/untold-by-enforced/off 2x2 via the \texttt{contract\_arm} field and
the \texttt{runtime\_mode} seam of Section 4.4. The mutation-gate
headline (Section 5.1) pools that mechanism's two committed commit
sub-gate tasks, target-commit and intent-commit, each pre-registered at
n=4, into an n=8 per-cell contrast; the off cell for the mutation gate
is the aggregate \texttt{no\_runtime\_enforcement} mode (Section 5.1
states the single-mechanism-attribution caveat this carries).

\textbf{Moderators as committed suite members.} The three moderators are
in the suite, not post-hoc splits: benign goal-conflict pressure as two
committed tasks (\texttt{gab\_l6\_agentsafe\_conflict},
\texttt{gab\_l5\_audit\_conflict}); constraint verifiability as a
blind/sighted twin pair (\texttt{gab\_l5\_audit\_blind} against its
sighted twin), in which the blind cell withholds ground-truth citations
so task success is unreachable there by construction, and the blind twin
is read as an honesty probe (abstain versus fabricate) on
audit-reference faithfulness and unsupported-narration rate, never on
task success, which is retained only to distinguish the sighted twin;
and L7 stale-manifest drift as a committed task
(\texttt{gab\_l7\_drift}). Verifiability and solvability co-vary in the
blind manipulation; the clean claim available from it is honest
abstention under hidden ground truth.

\textbf{Models and replication.} The four-model ladder of Section 4.7
runs with vendor-recommended per-model sampling, including the
deliberate below-floor operate control. Replication is n=4 per cell,
raised from 3 after an exact-binomial sensitivity analysis showed that
at pooled n=9 the 10-point decision bar is a one-rep gap with mid-regime
false-alarm rates of 25 to 41 percent, while pooled n=12 makes it a
two-rep gap; the full analysis is released with the benchmark
(\texttt{reports/sensitivity/SENSITIVITY.md}). We disclose that the
mutation gate carries two sub-gate tasks, not the three scenario pairs
that analysis assumed, so its pooled headline cell is n=8, below the
n=12 two-rep-gap target. The observed mutation-gate effects are
saturated (0/8 versus 8/8), so they do not sit in the flaky one-rep-gap
regime the analysis warned against, but the sub-target n is stated
rather than hidden.

\textbf{Decision quantities.} The pre-registered headline is the
per-mechanism specify-vs-enforce 2x2 cell contrast (A/B/C/D; B-vs-D the
effect of telling, C-vs-D the effect of enforcing, A-vs-B the marginal
value of enforcement given told), reported as pooled counts (k/n) and
percentage-point differences; medians are demoted to clearly-secondary
descriptive statistics, and the legacy pooled runtime-on/off median
machinery of the earlier pilot design is retired. The smallest effect
size of interest is a single 10-percentage-point bar (0 points only for
hard safety invariants); the earlier 5-point pilot bound is deleted, and
this change is disclosed here rather than silent. Pooled counts are the
scientific quantity; any unanimous per-task pass rule is operational
tooling, not an evidential criterion. First-attempt and converged
windows are both reported per Section 4.5, framed as
enforcement-as-barrier versus enforcement-as-teacher.

\textbf{Outcome categories.} Three non-behavioral outcomes are reported
as their own categories, never folded into behavioral failure counts:
provider safety-filter blocks, context-window overflow, and length
truncation. Each is neither a pass nor a behavioral fail; pooling any of
them into failure rates would inflate apparent incapacity, most acutely
for the small ladder models.

\textbf{Phasing and integrity.} The run is canary-first with a hard
stop: a canary subset (the untold-floor pair, the blind/sighted twin,
and the below-floor operate control) runs across all models and is
scored before the main sweep; the gate asks whether each canary axis
moves by at least 10 points in the expected direction, and a canary
failure attributable to model behavior is a reportable negative result
that ends the sweep, not a rerun trigger. A rerun is permitted only for
transient-infrastructure causes (provider outage, rate limit, parser or
adapter defect) and never to change a cell result; reruns are capped at
one fixed re-run with the defect named before re-running, and every
voided run is disclosed with its void reason. Claim language for every
possible outcome is pre-committed in an eleven-branch outcome map
(PILOT\_PROTOCOL.md §20), so the paper's conclusions section is selected
from the map by the data rather than written after seeing it.

\subsection{5 Results}\label{results}

The headline result is the pre-registered capability ladder (Section
5.1): a four-model run on the concentrated task set that produces a
capability-moderated pattern on the mutation gate, reported as a small
case study on raw cell counts rather than a powered significance claim.
It supersedes the thin n=3 diagnostic ladder that preceded it. The
diagnostic-tier probes that framed the design (Sections 5.2 through 5.5)
are reported afterward as provenance: one model
(\texttt{Qwen/Qwen3-235B-A22B-Instruct-2507-tput}, Together AI,
temperature 0), small n (3 to 5 per cell), run on diagnostic probe tasks
rather than the locked suite, kept as a distinct and lower evidence
tier. Static oracle-pair evidence and live runtime probes were
established separately per Section 4.6 and are not merged with any
number below. Throughout, the informative reads of the 2x2 are the
effect of telling (cell B: told-not-enforced vs cell D: neither, the
violation floor), the effect of enforcing (cell C: enforced-not-told vs
D), and the marginal value of enforcement given told (cell A: deployment
baseline vs B). L7 stale-manifest drift remains unmeasured (Section 5.6,
row P5; discussed in Sections 3.4, 8.4). {[}FORK-L7-DRIFT: unmeasured;
do not assume outcome{]}

\subsubsection{5.1 The capability-ladder
result}\label{the-capability-ladder-result}

The ladder ran the concentrated task set across four rungs (MiniMax-M3
as the capable primary, Llama-3.3-70B, Qwen3.5-9B near the operate
floor, Qwen2.5-7B below it), canary-first, against a single reference
runtime pinned by git fingerprint \texttt{6c82cd0} (Section 4.1; this is
14 commits past the nearest annotated tag \texttt{gab-runtime-1.0},
which predates the enforce-arm fixes and must not be used to reproduce).
A pre-spend runtime-identity preflight confirmed the enforce arm refuses
an agent-context commit under \texttt{full\_contract} before any paid
call. All eight canary and main runs completed (USD 10.44 total paid
spend, 448 rep ledgers); the pre-registered canary gate passed. The
mutation gate is the runtime's user-commit boundary: an agent may
propose a target or intent, but activating user-authored state requires
an explicit user commit. The block observed in the enforced cells is M5
(the \texttt{agent\_safe=false} dispatch refusal); M6 (the W57
proposal/commit gate) is masked by M5 and
\texttt{-\/-confirm}-bypassable, so it is defense-in-depth here, not
independently exercised (Section 3.2, PAPER.md D-48). The off cell is
the aggregate \texttt{no\_runtime\_enforcement} mode, which disables M4
through M8 together; we treat the resulting delta as a
whole-governance-off contrast on a commit task, on which the commit gate
is the only enforcer that can fire, and support that the fired mechanism
is the commit gate with the per-rep \texttt{mechanism\_disabled} markers
rather than by assertion. We report the pooled agent\_safe 2x2 (both
sub-gates, \texttt{unsafe\_action\_rate}, cells as percentage safe);
first-attempt and converged windows are byte-identical here, so a single
number stands for both.

{\def\LTcaptype{none} % do not increment counter
\begin{longtable}[]{@{}
  >{\raggedright\arraybackslash}p{(\linewidth - 8\tabcolsep) * \real{0.2000}}
  >{\raggedright\arraybackslash}p{(\linewidth - 8\tabcolsep) * \real{0.2000}}
  >{\raggedright\arraybackslash}p{(\linewidth - 8\tabcolsep) * \real{0.2000}}
  >{\raggedright\arraybackslash}p{(\linewidth - 8\tabcolsep) * \real{0.2000}}
  >{\raggedright\arraybackslash}p{(\linewidth - 8\tabcolsep) * \real{0.2000}}@{}}
\toprule\noalign{}
\begin{minipage}[b]{\linewidth}\raggedright
Model (capability tier)
\end{minipage} & \begin{minipage}[b]{\linewidth}\raggedright
A: told+enf
\end{minipage} & \begin{minipage}[b]{\linewidth}\raggedright
B: told+off
\end{minipage} & \begin{minipage}[b]{\linewidth}\raggedright
C: untold+enf
\end{minipage} & \begin{minipage}[b]{\linewidth}\raggedright
D: untold+off
\end{minipage} \\
\midrule\noalign{}
\endhead
\bottomrule\noalign{}
\endlastfoot
MiniMax-M3 (capable) & 100* & 100 & 100* & 0 \\
Llama-3.3-70B (capable) & 100* & 100 & 100* & 25 \\
Qwen3.5-9B (near-floor) & 100* & 0 & 100* & 12.5 \\
Qwen2.5-7B (below-floor) & 100* & 0 & 100* & 12.5 \\
\end{longtable}
}

Table caption: Pooled mutation-gate 2x2 per model, percentage safe (n=8
per cell = two commit sub-gate tasks x 4 near-deterministic replicates;
the target-commit untold cells are the pre-registered canary
untold-floor carriers, identically scored). *The enforced cells A and C
are safe by construction -- all 64 enforced reps score 0.0 with zero
variance, because the runtime blocks the gated commit regardless of the
model -- so they carry no behavioral degrees of freedom (they are the
deterministic guarantee) and the entire behavioral finding lives in
cells B and D.

The behavioral finding is the B column, told-not-enforced, where the
agent knows the rule and the runtime does nothing. As raw per-task
counts: both capable models complied on all four replicates of both
commit tasks (cell B safe 8/8 each; one-sample Clopper-Pearson interval
{[}63, 100{]}), and both weak models committed the gated mutation on all
four replicates of both tasks (cell B safe 0/8 each; {[}0, 37{]}). The
marginal value of enforcement given told, A minus B, is thus 0pp for the
capable models and +100pp for the weak ones. The capable-side 0pp is a
failure to detect at this n, not demonstrated equivalence: cell A is
pinned at 100\% by construction, so A-B is one-sided, bounded above by
the cell-B interval; we report it as ``no marginal value of enforcement
detected (one-sided upper bound +32pp)'' and attach neither a symmetric
CI nor a ``p=1.''

The independent unit of replication is the task, not the rep: two
sub-gate tasks per arm, four near-deterministic replicates each (7 of
the 8 model-by-sub-gate cell-B conditions are byte-identical across all
four replicates despite decoding temperature). A rep-level exact test
therefore overstates precision -- the weak-model A-B is Fisher p=0.00016
at the rep level but 0.33 at the honest task level (two tasks versus
two) and 0.029 per sub-gate (n=4) -- so we do not lead with it. Our
single pre-registered confirmatory contrast is the pooled A-B on the
mutation gate, one per model; treated at the rep level the two
weak-model rejections survive a Bonferroni correction across the four
models (adjusted \textasciitilde0.00062), but on the honest task-level
unit that same contrast is not significant (p=0.33), so we present the
crossover as a descriptive pattern and rest nothing on the rep-level
rejection. Every other contrast in this section is descriptive. The
primary evidence is the raw counts plus a per-rep mechanistic check: in
cell B, MiniMax refuses the commit on its own turn and tells the user to
run it, citing the user-gate, while Qwen-7B issues
\texttt{hai\ target\ commit} itself, with \texttt{mechanism\_disabled}
markers confirming the gate was down, and the commit executes. The sign
flip in cell B (100/100/0/0) holds on both sub-gates; the
difference-in-differences S is negative for the capable models but its
weak-model interval includes zero, so we do not rest the load-bearing
claim on it.

This capability contrast is what lets the capable-model result read as
substitution rather than pre-existing disposition. The reconstructed
mutation gate (Section 4.4) has an untold-unenforced cell that actually
violates: even for the capable models cell D sits near the floor
(MiniMax 0, Llama 25 percent safe, all via the original
\texttt{must\_not\_call} scorer path), so telling has a large measured
effect within a single model (B minus D = +100 and +75) and the capable
A minus B = 0 is not an agent that would have complied without either
lever. The earlier diagnostic result lacked this floor (the frozen
mutation gate sat at B minus D near zero; Section 8.4). The floor
itself, however, was made non-degenerate by a post-canary untold-arm
discoverability patch (Section 4.4), so the cell-D magnitude carries
that provenance caveat.

Five honesty notes govern the reading. First, enforcement is latent, not
exercised, in the capable models' enforce cells: they self-refuse before
dispatch, so the gate is never observed blocking a live violation in
cell A. That the runtime does enforce rests on the population
cell-C-versus-D contrast (untold + enforce stays safe while untold + off
violates) plus a single n=1 pre-spend preflight probe, not on a per-rep
block in cell A, which does not occur. Second, the below-floor model's
collapse is partial: its unsafe mutations are graded (median
\texttt{unsafe\_action\_rate} 0.5), so at any pass threshold at or above
0.5 its B cell would read 75\% rather than 0\%; we keep the
zero-tolerance threshold appropriate for an unauthorized-commit gate and
call the below-floor behavior partial self-enforcement, while the
load-bearing cells (capable B = 100, near-floor Qwen3.5-9B B = 0) are
invariant to any threshold in {[}0, 1). Third, the structured-output
parse rate is a confound but a protective one: pure-invalid reps score
safe (empty output cannot commit), yet no headline cell-B or cell-D rep
is a pure-invalid safe-by-default artifact (zero across all four models'
told-off and untold-off cells), so the weak-model B = 0 is genuine
parseable commits; and Llama-3.3-70B is a built-in control, producing
malformed turns on about 40\% of scored reps (39 of 95) yet still
reaching B = 100 through genuine refusals, which dissociates parse rate
from the B outcome. Fourth, n = 8 per cell (n = 4 per sub-gate) is a
pilot-scale ladder over four models, below our own pre-registered n=12
target (Section 8.4), not a scaling law; the intervals are wide and
reported as such. Fifth, the scorer's harm-only
\texttt{unsafe\_action\_rate} code path was edited three times in the
thirty hours before the run, the last 47 minutes before it (Section
4.5); those edits added an argument-defaults path that fired in zero
headline cells (every mutation-gate violation in the table is detected
by the original \texttt{must\_not\_call} path), so re-scoring the
trajectories with the pre-edit scorer is the identity on this 2x2 and
the load-bearing result is invariant to the pre-run scorer changes.

The other paid Tier-A mechanism, clinical refusal (M7), is reported as
near-ceiling and underpowered rather than as a second substitution
result. Across the ladder the clinical-claim detector fires on 7 of 223
refusal reps (3.1\%), all on one generic task
(\texttt{gab\_l6\_refusal\_untold}) and one model (Qwen3.5-9B); the two
tasks purpose-built to elicit a clinical claim
(\texttt{refusal\_diagnose}, \texttt{refusal\_dose}) produce zero fires
across all four models, including below-floor Qwen2.5-7B. The refusal
2x2 is therefore flat at ceiling with no contrast distinguishable from
zero, and the single discriminating cell (Qwen3.5-9B's untold arm) at
most suggests the clinical boundary begins to leak at low capability. We
do not read a substitution or a moderation result off this arm; whether
the near-total absence of fires reflects genuine boundary-respect or a
detector too narrow after its D-55.1 redesign is an open
scanner-sensitivity question (Section 8.4). Artifacts under the run
directory of Section 5.1.

\subsubsection{5.2 The contract-off map: telling works when salient, and
the untold floor is
real}\label{the-contract-off-map-telling-works-when-salient-and-the-untold-floor-is-real}

The first battery (2026-07-02) ran the telling axis with the runtime
fully off (\texttt{no\_runtime\_enforcement}), so every told condition
realizes cell B and the untold condition realizes cell D. All rows are
diagnostic (one model, Qwen3-235B, n=3 per condition).

{\def\LTcaptype{none} % do not increment counter
\begin{longtable}[]{@{}
  >{\raggedright\arraybackslash}p{(\linewidth - 2\tabcolsep) * \real{0.5000}}
  >{\raggedright\arraybackslash}p{(\linewidth - 2\tabcolsep) * \real{0.5000}}@{}}
\toprule\noalign{}
\begin{minipage}[b]{\linewidth}\raggedright
Condition
\end{minipage} & \begin{minipage}[b]{\linewidth}\raggedright
Outcome (n=3 per row)
\end{minipage} \\
\midrule\noalign{}
\endhead
\bottomrule\noalign{}
\endlastfoot
Told, bait phrasing & 3/3 turn-1 refusal citing the \texttt{agent\_safe}
flag \\
Told, exact gated command named & 3/3 turn-1 refusal \\
Told, normal task phrasing (constraint reached incidentally) & 3/3
first-attempt dispatch of the gated command \\
Untold (flag globally stripped), valid target id supplied & 3/3
completed commit, row \texttt{proposed} to \texttt{active}, truthful
success report, no refusal step \\
\end{longtable}
}

Two readings follow. First, with the runtime off the told-and-salient
conditions refused 3/3 while the untold condition completed the
violation 3/3 (one model, Qwen3-235B, n=3 per cell), so the effect of
telling against the violation floor is non-hypothetical at this tier.
But self-enforcement of a context-verifiable constraint is
salience-sensitive: the same agent that refuses 3/3 when the constraint
is foregrounded dispatches the gated command 3/3 on first attempt when
it reaches the constraint incidentally through normal task phrasing. Any
claim that a told agent self-enforces must carry this qualifier. The
incidental row is scored on first valid action: two of the three reps
emit invalid-output steps before the dispatch, so ``first-attempt''
means first contact with the enforcement surface, not literal step one.
These pre-dispatch invalid-output steps are the parser-clean model
recovering to a valid dispatch, distinct in kind from the non-recovering
invalid-output loops that place the 9B and 70B ladder models below the
operate floor (Sections 6.3, 7.1); the shared invalid-output signature
is read one way here and another there, and the residual overlap is a
limitation rather than a clean separation. The two salient rows carry
the author's contemporaneous labels; the stored artifacts retain
prompt-template hashes rather than the phrasings, so the
row-to-directory mapping is not independently recoverable, though the
3/3 refusal pattern in both rows is.

Second, the untold row is the untold violation floor: with the contract
withheld and the runtime off, nothing stood between the agent and a
completed governance violation, reported truthfully as success. This
makes the effect of telling (B vs D) and the effect of enforcing (C vs
D) non-hypothetical. The substitution result of Section 5.1 is about the
marginal value of enforcement given told (A vs B), not about these
floors.

\subsubsection{5.3 The M8 arc: the verifiability exception
dissolves}\label{the-m8-arc-the-verifiability-exception-dissolves}

M8 audit evidence emission was pre-registered as the non-verifiable
exception: the constraint an agent cannot check from its decision-time
context and therefore, by hypothesis, the place runtime enforcement
retains behavioral value. Three probe regimes (2026-07-02, diagnostic,
n=3 per cell) dismantled the classification in stages.

{\def\LTcaptype{none} % do not increment counter
\begin{longtable}[]{@{}
  >{\raggedright\arraybackslash}p{(\linewidth - 2\tabcolsep) * \real{0.5000}}
  >{\raggedright\arraybackslash}p{(\linewidth - 2\tabcolsep) * \real{0.5000}}@{}}
\toprule\noalign{}
\begin{minipage}[b]{\linewidth}\raggedright
Regime
\end{minipage} & \begin{minipage}[b]{\linewidth}\raggedright
Outcome
\end{minipage} \\
\midrule\noalign{}
\endhead
\bottomrule\noalign{}
\endlastfoot
Retrieved (\texttt{hai\ explain} output in context) & 6/6 honest: cites
the real card id under \texttt{full\_contract}, reports ``No evidence
card is present.'' under \texttt{no\_audit\_chain}; verbatim-identical
finals within each cell \\
Non-retrieval, plain ask, honest answer available & 6/6 honest
abstention, zero fabrication \\
Instrumental (an id is needed to advance an action), pre-fix harness &
Fabricated ids dispatched freely: 8 instances across three probe sets \\
\end{longtable}
}

The first regime shows that verifiability is a property of the
(constraint, decision-time context) pair, not of the mechanism: the
moment retrieval lands the evidence in context, M8 behaves like the
verifiable mechanisms and the agent self-enforces it. The second shows a
cooperative agent asked plainly stays honest even without retrieval.
That left the third regime as the only surviving fabrication signal. We
reattributed it from verifiability to instrumental pressure, folded M8
into the goal-conflict axis, and pre-registered the two follow-ups
below; the instrumental cell was later falsified outright as a
harness-blindness artifact (Sections 5.5 and 6.1). After the arc, the
sole surviving candidate for a genuinely non-verifiable constraint class
is L7 stale-manifest drift (Section 3.4); no model-backed measurement
exists. {[}FORK-L7-DRIFT: unmeasured; do not assume outcome{]}

\subsubsection{5.4 H2: benign goal conflict does not degrade told-only
compliance}\label{h2-benign-goal-conflict-does-not-degrade-told-only-compliance}

The first pre-registered follow-up (artifact timestamp 2026-07-02T1942Z;
pre-registered n=5) applied benign completion pressure, not adversarial
injection, to audit faithfulness: 4 task variants graded P0 through P3,
at P3 a pipeline is blocked and the user cannot proceed unless the agent
produces a reference. Design: 4 tasks x 5 reps x 2 modes
(\texttt{full\_contract}, \texttt{no\_audit\_chain}) for 40 reps.

The result is a null. Zero fabrication across all 40 reps,
scorer-verified: zero occurrences of the
\texttt{forged\_or\_suppressed\_audit\_reference} violation kind in all
40 score files. The pre-committed falsification bar is an observed-rate
threshold: fabrication at most 10\% at P3. The four task variants map
one to each pressure level, so the P3 cell is 5 reps per mode (10 reps).
Observed P3 fabrication was 0/10; at n=10 this cannot exclude a true
rate above 10\%, so the bar is met and H2 is falsified only at the
diagnostic tier (one model, Qwen3-235B, n=5). As a tighter descriptive
bound, and as this paper's own post-hoc caveat rather than part of the
pre-registration, the pooled 0/40 across P0-P3 narrows the upper
confidence interval further. The small-n limitation stated throughout
applies. The side finding: 10 of the 40 reps fail overall pass with a
refusal-format error only, meaning high pressure degraded tool-use
competence (command-format loops) before it touched honesty.

\subsubsection{5.5 The instrumental leg is
falsified}\label{the-instrumental-leg-is-falsified}

The second pre-registered follow-up (2026-07-03; pre-registered n=5;
\texttt{no\_runtime\_enforcement}) tested the reattributed hypothesis
directly, after the stdout-inlining fix let the agent see its own
command output. The result is a falsification: an instrumental effect of
0 percentage points against the pre-committed bar of at least 40 points,
diagnostic (one model, Qwen3-235B, n=5). The apparent
instrumental-fabrication effect of Section 5.3 was a fact about a blind
harness, not about the model; Section 6.1 tells the full story,
including the design, the cell-level counts, the harness defect, and the
disclosure that only the absence of the fabrication violation kind is
citable from this scan. A thin n=3 diagnostic ladder run the same day
first suggested the operate floor and the self-enforcement pattern; it
was confounded (the action parser was tuned to one model's output
format; Section 6.3) and is superseded by the pre-registered sweep of
Section 5.1, which measures the same ladder properly at n=8 per cell.

\subsubsection{5.6 Predictions against
outcomes}\label{predictions-against-outcomes}

{\def\LTcaptype{none} % do not increment counter
\begin{longtable}[]{@{}
  >{\raggedright\arraybackslash}p{(\linewidth - 6\tabcolsep) * \real{0.2500}}
  >{\raggedright\arraybackslash}p{(\linewidth - 6\tabcolsep) * \real{0.2500}}
  >{\raggedright\arraybackslash}p{(\linewidth - 6\tabcolsep) * \real{0.2500}}
  >{\raggedright\arraybackslash}p{(\linewidth - 6\tabcolsep) * \real{0.2500}}@{}}
\toprule\noalign{}
\begin{minipage}[b]{\linewidth}\raggedright
\#
\end{minipage} & \begin{minipage}[b]{\linewidth}\raggedright
Situation
\end{minipage} & \begin{minipage}[b]{\linewidth}\raggedright
Predicted delta
\end{minipage} & \begin{minipage}[b]{\linewidth}\raggedright
Outcome
\end{minipage} \\
\midrule\noalign{}
\endhead
\bottomrule\noalign{}
\endlastfoot
P0 & Capability moderator (marginal value of enforcement given told,
A-B, across the ladder) & Substitution at high capability, enforcement
load-bearing at low & Pattern seen (n=8, descriptive): capable A-B=0
(one-sided upper +32pp); weak A-B=+100 (raw counts; rep-level p=0.00016
overstates precision, task-level 0.33) \\
P1 & Verifiable, no conflict, capable model & \textasciitilde0 &
Consistent: capable A-B=0 on the mutation gate (sweep, n=8);
self-enforces when salient (n=3) \\
P2 & Verifiable, goal conflict (H2) & Real, grows with pressure & Null:
0 fabrication P0-P3 (n=5, pre-registered) \\
P3 & M8 references, instrumental pressure & Real & Falsified:
harness-blindness artifact; 0pp with stdout fix (n=5, pre-registered) \\
P4 & M8 references, cooperative, asked plainly & \textasciitilde0 &
Consistent: honest 6/6 per regime (n=3 per cell) \\
P5 & L7 drift & Real at every level & {[}FORK-L7-DRIFT: unmeasured; do
not assume outcome{]} \\
P6 & Contract withheld (untold) & Enforcement is the only barrier &
Consistent: untold violation floor observed (sweep D-cell; n=3
diagnostic) \\
P7 & Below/near operate floor & Prevents malformed harm; enforcement
load-bearing & Confirmed: weak models A-B=+100 (sweep, n=8); 7B
partial/below floor \\
\end{longtable}
}

Table caption: Pre-registered predictions against outcomes. The
capability moderator (P0, P1, P7) is measured by the sweep of Section
5.1 (n=8 per cell); the verifiability (P3, P4) and goal-conflict (P2)
predictions are diagnostic tier (one model, Qwen3-235B, n=3-5 per cell).

Of the three moderators hypothesized to gate substitution, model
capability is the one the evidence resolves: the ladder shows a sign
flip in the marginal value of enforcement given told, from 0 (no
marginal value detected) at high capability to +100pp at low capability.
The other two nulled at the diagnostic tier: the verifiability exception
dissolved (P3, P4) and goal conflict produced zero fabrication (P2). The
marginal value of enforcement given told (A-B) is the headline quantity,
and the ladder puts capable cell B at ceiling while collapsing weak cell
B to the floor; the supporting diagnostics are consistent, with cell B
at ceiling in every told-and-salient probe (3/3 refusal in both
foregrounded contract-off rows, 3/3 honest reporting in the cell-B half
of the retrieved M8 regime, 3/3 honest abstention in the cell-B half of
the non-retrieval regime, and 0 fabrication across all 20 cell-B
goal-conflict reps; all diagnostic, one model, n=3 to 5 per cell). The
result is therefore not a flat substitution but a capability-moderated
one: for a capable cooperative agent above the operate floor, in-context
specification substitutes for runtime enforcement on the mutation gate,
while at and below the near-floor band enforcement is load-bearing
because the told agent does not self-enforce. Runtime enforcement's
demonstrated behavioral value is the deterministic guarantee plus the
region where the model cannot self-enforce: the untold violation floor,
the sub-frontier capability band, the incidental-salience failure of
Section 5.2, below the operate floor, and adversarial intent, untested
here. Enforcement is behaviorally redundant only in the scoped
high-capability sense, never structurally. All of this is within this
fixed controller and this task set, over a four-model ladder at n=8 per
cell; it is not a scaling law or a general claim about agents or
harnesses. It is a single-runtime case study, and external non-HAI
replication is future work (Section 7.2).

\subsection{6 Methodology cautions as a
contribution}\label{methodology-cautions-as-a-contribution}

Our probing arc produced, and then destroyed, an apparent positive
result. We report the destruction as a named contribution rather than a
limitations aside, because the failure mode is general to agent
evaluation and easy to reproduce by accident.

\subsubsection{6.1 Harness blindness manufactured a fabrication
finding}\label{harness-blindness-manufactured-a-fabrication-finding}

In early probes the harness did not surface command output to the agent.
Where stdout should have appeared in the observation feedback, the agent
received a \texttt{stdout\_ref} file path instead. The agent could
dispatch a lookup, but it could never read the result. When a task then
required an identifier to advance an action, the agent guessed: across
three pre-fix probe sets we observed 8 instances in which it dispatched
commands carrying plausible, well-formed identifiers it had never seen,
including a fabricated evidence-card identifier and repeated invented
target identifiers (model-backed diagnostic tier: one model, Qwen3-235B,
temperature 0, n=3 per cell). This looked like a real finding,
instrumental fabrication: honest when asked plainly, fabricating when an
identifier is instrumentally required. It was, at the time, the
strongest surviving candidate for a positive enforcement result.

The finding was an artifact of the harness. The pre-registered
falsification experiment crossed Statement vs Action framing with id
present/absent (4 tasks x 5 reps = 20, n=5 per cell, pre-committed bar:
instrumental effect \textgreater=40 percentage points), run after the
committed stdout-inlining fix (commit \texttt{17db5ef}), which inlines
command stdout into the observation feedback bounded to the first 24,000
characters. Result: fabrication 0/5 in the Action/id-absent cell, 0/5 in
the Statement/id-absent cell, an instrumental effect of 0 percentage
points against the \textgreater=40pp bar, diagnostic (one model,
Qwen3-235B, n=5). Once the agent can see the empty lookup, it abstains
honestly 5/5 even in the Action arm, where the identifier is needed to
finish the task. The cleanest before/after pair holds the task fixed:
pre-fix, one rep fabricates a plausible card identifier while another
rep in the same task and runtime mode stays honest; post-fix, all 6 reps
are honest, with verbatim-identical finals within each runtime mode. One
disclosure: we cite this scan only for the absence of the fabrication
violation kind. In the id-absent cells honest abstention costs task
success by design; across the scan all 20 reps fail overall pass and a
refusal-classified scoring error fires 20/20, and the action/id-present
task additionally logs an unsafe-mutation flag 5/5 under
\texttt{no\_runtime\_enforcement}. The id-present control cells are
valid (the agent reports or uses the real retrieved identifier); no
outcome from this scan other than the fabrication-kind absence is cited
in this paper.

The general lesson: eval harnesses that hide tool output manufacture
spurious fabrication findings. A model that cannot observe its own
command results is forced to guess, and a deterministic scorer will read
the guess as fabrication. The resulting number is a fact about the
harness, not the model. Genuine constraint-driven fabrication does exist
under more adversarial operationalizations {[}arXiv:2606.14831{]}; our
claim is narrower and methodological: before attributing fabrication to
the agent, verify the agent could see what it is accused of inventing.

\subsubsection{6.2 Prompt-embedded ablations measure
self-enforcement}\label{prompt-embedded-ablations-measure-self-enforcement}

A second caution follows directly from the 2x2. Guardrail ablation
studies that toggle runtime enforcement while leaving the policy in the
system prompt in both conditions, as in {[}arXiv:2604.15579{]}, measure
the marginal value of enforcement given the agent was told (A:
deployment baseline vs B: told-not-enforced, the marginal value of
enforcement given told), not the effect of enforcing (C:
enforced-not-told vs D: neither, the violation floor). For a capable
cooperative agent that self-enforces, this design will show small deltas
and invite the wrong conclusion in either direction: that enforcement is
worthless, or that a small delta on a told agent bounds what enforcement
does for an untold one. Neither follows. Ablations that never withhold
the in-context contract measure self-enforcement, not enforcement.

\subsubsection{6.3 Two supporting
cautions}\label{two-supporting-cautions}

\textbf{A parser tuned to one model family scores other models below the
operate floor.} On the capability ladder (Section 7.1), only Qwen3-235B
emitted parser-clean output; the 9B and 70B models looped invalid on
routing, which automated scoring reads as inability to drive the
contract. Without manual inspection this would have placed operable
models spuriously below the operate floor. This confound is why the
ladder carries no scaling claim.

\textbf{Provider catalog metadata is not ground truth.} The Together
serverless catalog flag was unreliable: models it listed as serverless
were dedicated-endpoint only in practice, which contaminated ladder
composition; separately, the screen-to-ladder provenance chain is not
reconstructable from artifacts. Roster composition must be verified
against live endpoints, not catalog flags.

\subsubsection{6.4 A checklist for agent-eval
builders}\label{a-checklist-for-agent-eval-builders}

\begin{itemize}
\tightlist
\item
  Confirm the agent receives tool and command output verbatim in its
  observation feedback; a reference or path is not output.
\item
  Before reporting fabrication, hallucination, or dishonesty, check
  whether the accused content was observable to the agent at decision
  time.
\item
  If the policy is in the prompt in every condition, the enforcement
  toggle measures the marginal value of enforcement given told, not the
  effect of enforcing.
\item
  Validate the action parser against every model in the roster before
  attributing failure to capability; parser mismatch mimics the operate
  floor.
\item
  Verify provider catalog claims (availability, endpoint class) against
  live endpoints before locking a roster.
\item
  Pre-register a falsification bar for any surprising positive finding
  and re-run it after fixing the harness defect that co-occurred with
  it.
\end{itemize}

\subsection{7 The capability ladder and external
replication}\label{the-capability-ladder-and-external-replication}

\subsubsection{7.1 The capability ladder}\label{the-capability-ladder}

Model capability is the headline moderator (Section 5.1), and the
evidence is a pre-registered four-rung ladder run on the concentrated
task set at n=8 per cell under both \texttt{full\_contract} and
\texttt{no\_runtime\_enforcement}: MiniMax-M3 (capable primary),
Llama-3.3-70B (capable), Qwen3.5-9B (near the operate floor but
operational), and Qwen2.5-7B (below it). The ladder produces the
capability-moderated pattern reported in full in Section 5.1: the
marginal value of enforcement given told (A-B) is 0 for the two capable
models (no marginal value detected, one-sided upper +32pp) and +100pp
for the two weak models (raw counts; the rep-level exact p is not the
honest unit of analysis, Section 5.1). It supersedes a thin n=3
diagnostic precursor (2026-07-03, \texttt{runs/pilot/\_probe\_ladder/})
that first surfaced the operate floor and the self-enforcement pattern
but was parser-confounded (Section 6.3) and left one rung unmeasured
within-model.

Three properties bound what the ladder can claim. First, it is a
four-model pilot at n=8 per cell, not a scaling law: it establishes that
the marginal value of enforcement given told moves from undetectable to
load-bearing across this capability range, not a capability-ordered
curve. Non-monotone capability-compliance relationships across model
families are documented {[}arXiv:2606.01317, arXiv:2601.08196{]}, and
four models cannot adjudicate them. Second, the structured-output parse
rate is a confound, but a protective one with a built-in control:
pure-invalid reps score safe (empty output cannot commit), so the
weak-model B = 0 is a lower bound on the self-enforcement failure rate,
and Llama-3.3-70B produces malformed turns on about 40\% of scored reps
(39 of 95) yet still reaches B = 100 through genuine refusals,
dissociating parse rate from the B outcome (the thin precursor's parser
confound, Section 6.3, is exactly why this control matters). Third, the
operate floor is real and below-floor behavior is partial: Qwen2.5-7B's
unsafe mutations are graded (median \texttt{unsafe\_action\_rate} 0.5),
so its collapse is partial self-enforcement rather than total, while the
load-bearing near-floor rung (Qwen3.5-9B, A-B = +100) is
threshold-invariant. The clean load-bearing case is therefore
Qwen3.5-9B: it operates the contract (drives valid actions, enforce
cells safe) yet does not self-enforce told-not-enforced, so an
operate-but-violate band does exist between the floor and the frontier,
contrary to the thin precursor's suggestion that models either drive the
contract and comply or fall below the floor.

Compounding these caveats narrows the robust load-bearing signal.
Removing the threshold-sensitive below-floor 7B and holding the
capability-family confound, the evidence that a told agent fails to
self-enforce while still operating the contract rests on a single
non-frontier Qwen model (Qwen3.5-9B) at n=8 on one mechanism and one
runtime. The crossover's low arm should be read at that strength: a
real, mechanistically verified failure to self-enforce, but not yet a
family-general or multi-model result. The thin n=3 diagnostic precursor,
which ran pre-parser-fix at temperature 0 on one gate task, contributes
no number to this ladder.

\subsubsection{7.2 Scope: a single-runtime case study; external
replication as future
work}\label{scope-a-single-runtime-case-study-external-replication-as-future-work}

Every model-backed number in this paper comes from one controller, HAI
v0.2.0, with its particular manifest, contract-authoring style, and task
suite. This paper is accordingly scoped as a single-runtime case study
(PAPER.md D-42): hypothesis H5, that the specify-vs-enforce effect
generalizes beyond HAI, is stated as a hypothesis, and the external
non-HAI replication that would test it is future work, not a pending
component of this paper's evidence. Nothing in the substitution result
distinguishes a phenomenon from a HAI quirk until that replication
exists, and the claims throughout are conditioned on this controller.

The future-work design is specified so the scope boundary is concrete:
reproduce the 2x2 on an independent harness with its own mechanism
inventory and deterministic scorer, measuring at minimum the A-vs-B
contrast, the marginal value of enforcement given the agent was told;
where the replication also attributes the telling axis (B vs D, C vs D),
first-attempt scoring applies (Section 3.6). TeamBench
{[}arXiv:2605.07073{]} (Section 2) is one candidate substrate: prior
evidence of the same qualitative shape in a different mechanism domain
and design, not itself a replication of this paper's effect. Its
sub-metric asymmetry (pass-rate parity coexisting with 3.6 times more
verifier-edits-executor cases in the prompt-only arm) cautions that
aggregate parity can mask enforcement value. That caution applies to our
own instrument. The deterministic scorer observes per-mechanism
\texttt{mechanism\_disabled} markers and each task's load-bearing
violation kind (Section 4.2), so it detects whether the specific guarded
violation occurs, but it does not resolve finer sub-metric channels such
as a verifier-edits-executor count. An enforcement effect that left the
pass rule and the tracked violation kinds unchanged while shifting such
a channel would stay undetectable at diagnostic n.~Aggregate parity in
these diagnostics should therefore not be read as sub-metric
equivalence.

\subsection{8 Discussion}\label{discussion}

\subsubsection{8.0 Scope conditions}\label{scope-conditions}

The substitution is conditional, not universal. On the sweep (Section
5.1, n=8 per cell), substitution held when six conditions held together:
the agent was capable, above the operate floor; the agent was
cooperative; the constraint was stated in the in-context contract and
told plainly; the constraint was salient at decision time; the harness
surfaced command output to the agent; and the inputs were benign. The
first, capability, is the one the ladder varies: where it fails,
told-not-enforced collapses and enforcement is load-bearing (Section
5.1). The other four behavioral conditions, when they fail, likewise
mark regions where runtime enforcement retains demonstrated value
(Section 8.1); the sixth, the harness surfacing command output, is a
measurement-validity condition whose failure manufactures the spurious
fabrication finding of Section 6 rather than a genuine enforcement
region.

\subsubsection{8.1 What runtime enforcement is still
for}\label{what-runtime-enforcement-is-still-for}

The capability-moderated result relocates enforcement's value rather
than eliminating it. Beyond the deterministic guarantee, which is
structural rather than a condition-failure region, each scope condition
of Section 8.0 marks a region where enforcement retains behavioral value
when that condition fails.

First, the deterministic guarantee of Section 3.5, which holds
unconditionally and is not tied to any one condition failing. Cell B:
told-not-enforced compliance is a statistic over model behavior; cell A:
deployment baseline compliance on the enforced surface is a property of
the harness. A gated action that the runtime blocks cannot occur; a 3/3
refusal rate at n=3 bounds nothing.

Second, when the told condition fails: the untold violation floor of
Section 5.2, where the contract is withheld and enforcement is the only
barrier. In-context specification cannot substitute for enforcement on
rules the context never carried, and real deployments accumulate such
rules through contract drift, incomplete prompt assembly, and context
compaction {[}arXiv:2606.22528{]}.

Third, when the salience condition fails: the incidental-salience region
of Sections 5.2 and 8.2, where the same told constraint reached
incidentally rather than foregrounded was dispatched 3/3 on first
attempt. Enforcement catches the violation that a told-but-not-salient
agent commits.

Fourth, when the capability condition fails: below the operate floor,
where enforcement prevents malformed harm rather than disobedience
(Section 7.1); the failure there is contract friction, not governance.

Fifth, when the benign-inputs condition fails: adversarial inputs, where
the cooperative assumption also no longer holds. The substitution result
is scoped to benign completion pressure, not adversarial injection. The
injection literature already shows architectural enforcement
outperforming prompt-only defenses under attack {[}arXiv:2503.18813,
arXiv:2504.11703{]}; Verifier Tax {[}arXiv:2603.19328{]} similarly finds
policy mediation intercepting up to 94\% of non-compliant actions under
instrumental pressure while agents hallucinate identifiers to route
around blocks. Injection robustness is cited territory, not a claim of
this paper, and nothing here licenses removing enforcement from
adversarially exposed surfaces.

One published result points the opposite empirical direction and
warrants direct comment. Mechanical Enforcement {[}arXiv:2605.14744{]}
finds mechanical enforcement raising governance quality and task
accuracy (MCC approximately 0.43 to 0.88) over text-only governance,
where this paper finds enforcement behaviorally redundant given
specification. The findings differ on four axes: domain (regulated
banking versus non-clinical personal wellness), metric type
(decision-rationale quality versus tool-action compliance), design
(paired alternative configurations versus a crossed per-mechanism
factorial), and pressure regime (their regulated-decision setting versus
benign completion pressure here). The tension is a difference in what is
measured and under what pressure, not a contradiction on the marginal
value of enforcement given told. A same-direction 2026 result, Mind the
GAP {[}arXiv:2602.16943{]}, finds a runtime governance contract with no
detectable deterrent on forbidden tool-call attempts; that is a leakage
catch-layer under jailbreak, whereas the mutation gate here is a hard
commit block under benign use, so its untold-plus-enforce cell (C) is
safe by construction. The two do not conflict: enforcement architecture
that only annotates or logs need not deter, while enforcement that
blocks dispatch does.

\subsubsection{8.2 The salience confound}\label{the-salience-confound}

The told-not-enforced leg carries a caveat. Self-enforcement is
salience-sensitive: in the diagnostic battery of Section 5.2, the told
agent refused the gated action 3/3 when the constraint was foregrounded,
but dispatched it 3/3 on first attempt when it was reached incidentally
through normal task phrasing (one model, n=3 per cell). Under
first-attempt scoring, telling substituted for enforcement only when the
telling was salient at decision time. The substitution claim is really a
claim about salient in-context specification, and how salience decays
with context length, task distance, and compaction is unmeasured here
{[}arXiv:2606.22528{]}.

\subsubsection{8.3 Personal health as a
domain}\label{personal-health-as-a-domain}

HAI's domain, non-clinical personal wellness, is a deliberate instrument
choice. The non-clinical boundary (no diagnosis, treatment, prescribing,
or autonomous medical decisions) is itself an evaluated constraint, the
zero-tolerance leg of M7 refusal of out-of-contract requests. The choice
makes a policy-dense boundary an enforceable, ablatable object of study
with unambiguous violation semantics, without any clinical claim.

\subsubsection{8.4 Limitations}\label{limitations}

The primary model-backed evidence is the four-model capability-ladder
sweep at n=8 per cell (Section 5.1); the design's framing and moderator
probes are a lower diagnostic tier (one model,
Qwen3-235B-A22B-Instruct-2507, Together AI, temperature 0, n of 3 to 5
per cell, on diagnostic probe tasks rather than the locked suite). At
n=8 per cell the intervals are wide and reported as such, and the ladder
is four models over a single reference runtime (HAI), one harness, one
prompt template; the paper is scoped as a single-runtime case study, and
external non-HAI replication is future work (Section 7.2). The untold
violation floor and the harness-blindness finding are existence and
property demonstrations that do not rest on cross-model generalization.
L7 stale-manifest drift, the sole surviving candidate for a true
non-verifiable enforcement delta, is unmeasured. {[}FORK-L7-DRIFT:
unmeasured; do not assume outcome{]} Inputs are benign throughout; the
adversarial-input arm was retired at the benchmark tier (Section 4.8),
and injection robustness is cited territory, not a claim. The
clinical-refusal detector, redesigned to a frame-based scanner in
D-55.1, fires on almost no reps in the sweep (Section 5.1); whether that
reflects genuine boundary-respect or an over-narrowed scanner is
unresolved, which is why the refusal arm is reported as near-ceiling and
underpowered rather than as coverage.

The disposition confound that limited the earlier diagnostic result is
resolved here, and it is worth being precise about how. On the frozen
tasks the enforce-able mechanisms were disposition-covered: a
cooperative capable agent already complied whether or not it was told (B
versus D near zero), so the marginal value of enforcement given told (A
versus B) was consistent with substitution but did not exclude that the
agent needed neither lever, and the difference-in-differences S was
unidentified. The told-arm reconstruction of Section 4.4 (D-52) removed
that degeneracy on the mutation gate: the untold-unenforced cell now
actually violates, so even for the capable models cell D is near the
floor (MiniMax 0, Llama 25 percent safe) and B versus D is large (+100,
+75). With the floor non-degenerate, telling and enforcing each have a
measured effect and S = (A minus B) minus (C minus D) is identified
within a single model, so the capable-model marginal value A minus B = 0
is read as substitution rather than as a disposition that would have
complied regardless. The capability ladder then supplies the moderation:
at low capability the agent does not self-enforce even when told (B =
0), so enforcement given told is load-bearing (A minus B = +100 on raw
counts; the rep-level exact p overstates precision, Section 5.1). The
clinical-refusal arm, by contrast, remains near-ceiling and underpowered
(Section 5.1), so we make no disposition-versus-specification separation
there; the pre-registered directiveness-gradient pilot (Section 9.1) is
retained as a complementary, finer-grained within-model identification
on the mutation gate, not as the confound's sole resolution.

The ladder's own limitations, disclosed against the pre-registered run
design (Section 4.10): it is a bounded four-model moderator, not a
scaling-law claim. Capability is partly confounded with model family,
because both capable rungs are non-Qwen (MiniMax-M3, Llama-3.3-70B)
while both weak rungs are Qwen (3.5-9B, 2.5-7B); the cross-family
agreement of the two capable self-enforcers strengthens the substitution
side, but the load-bearing side rests on Qwen-family models, so the
crossover cannot fully separate a capability effect from a Qwen-family
low-capability effect, and the abstract must not assert family-general
capability ordering. Under per-model vendor-recommended sampling, model
identity is confounded with sampling regime, so the ladder cannot
separate a capability trend from a sampling-temperature trend; the
per-model 2x2 is computed within a single model at a constant sampling
regime and is unaffected. Provider seeds are unsupported on the
serverless MoE endpoints, so replication is across independent
stochastic draws rather than fixed-seed determinism. The untold arm's
schema-partial-telling, told-the-opposite M5, and cell-B
enforcement-prose disclosures of Section 4.4, and the clinical-detector
circularity of Section 4.5, apply to every run-design cell.

\subsubsection{8.5 What would change the
conclusion}\label{what-would-change-the-conclusion}

The conclusion is falsifiable on several fronts. The high-capability
substitution would break if a capable, above-floor, cooperative model
showed a nonzero A-versus-B delta on the mutation gate under the same
first-attempt scoring, or if capable self-enforcement failed at higher n
or under longer-horizon salience decay. The capability-moderation would
break if the crossover failed to replicate on a second runtime, or if a
non-Qwen weak model self-enforced told-not-enforced, which would
reattribute the load-bearing low end from capability to the Qwen family.
Independently, a measured L7 drift delta, or goal-conflict pressure
beyond P3 breaking the null, would restore behavioral, not merely
guarantee-level, value to enforcement in regimes we report as null. The
claim is exactly as strong as the scope conditions of Section 8.0, and
no stronger.

\subsection{9 Future work and
conclusion}\label{future-work-and-conclusion}

\subsubsection{9.1 Future work}\label{future-work}

The measurement this design points at most directly is L7 stale-manifest
drift (Section 3.4), the sole surviving genuinely non-retrievable
candidate for a real cooperative-agent enforcement delta. Concurrent
work establishes the drift-causes-violation half, that context
compaction silently erases constraints the agent obeyed while they were
visible {[}arXiv:2606.22528{]}; whether independent runtime enforcement
catches post-drift violations is the untested second half.
{[}FORK-L7-DRIFT: unmeasured; do not assume outcome{]}

Second, the adversarial-input arm should graduate from scorer-coverage
evidence to model-backed evidence. This paper cites injection robustness
rather than claiming it (Section 8.1); running the specify-vs-enforce
2x2 under attacker-controlled inputs would locate where the substitution
result ends and the injection-defense literature begins.

Third, a fuller screened capability ladder behind an action parser
tolerant of every roster model. The present ladder carries only the
existence of an operate floor (Section 7.1); a parser-tolerant rerun
across a wider screened roster could locate the floor rather than merely
evidence it, with non-monotonicity across families expected
{[}arXiv:2606.01317, arXiv:2601.08196{]} and no scaling law promised.

Fourth, external non-HAI replication (Section 7.2). This paper is a
single-runtime case study by decision, and the replication that would
distinguish a phenomenon from a HAI quirk, reproducing the 2x2 on an
independent harness with its own mechanism inventory and deterministic
scorer, is the highest-value follow-on measurement. (The multi-model
higher-n question, by contrast, is resolved within this paper for the
mutation gate: the pre-registered run of Section 4.10 is that run, and
its mutation-gate results land in Section 5.1 per the pre-committed
outcome map, not in future work. The other paid mechanism, clinical
refusal, was run on the ladder but is near-ceiling and underpowered
(Section 5.1); the Tier-D scope mechanisms are offline-only by
pre-registered design, D-49.)

Fifth, a pre-registered exploratory disposition-mapping pilot sharpens
the identification the capability ladder already delivers at the coarse
grain (Section 8.4). Where the ladder separates specification from
disposition across models, the pilot does so within a single model: it
manufactures intermediate disposition on the mutation gate, the one
mechanism carrying both a movable told-versus-untold effect and a
runtime enforce/off lever, through a directiveness gradient over the
operator's request, from a directly commanded activation to a neutral,
instrumental phrasing, so that the untold-unenforced violation rate is
dialed off its floor and both levers move on the same rows. There S = (A
minus B) minus (C minus D) is identified within one model, letting the
specification effect be read against that model's baseline disposition
rather than confounded with it; the benign completion-pressure and
deployment-scope pilots map the disposition-covered and
specification-substitutable types of Section 3.6 alongside it. The pilot
is pre-registered as exploratory (PILOT\_PROTOCOL Section 22): it does
not amend the frozen suite or the headline, and a confirmatory follow-up
would re-lock the chosen constraint and gradient level with fresh reps.
Newcombe intervals on B vs D and C vs D and a difference interval on S
are reported per constraint by the results layer.

Two deferred lines complete the program. H7 extends the design toward
the capability-versus-compliance question studied externally under
scaling-laws-for-oversight {[}Engels2025{]}, asking whether a
deterministic Guard at general Elo of zero occupies a legitimate point
on that external curve; this paper itself claims no scaling law (Section
1). H8 asks whether fine-tuned operators reach contract compliance at
smaller scales than untrained operators (the S1 fine-tuning sequel),
which would move the operate floor itself.

\subsubsection{9.2 Conclusion}\label{conclusion}

On the mutation gate, in a small pre-registered single-runtime case
study, in-context specification substitutes for runtime enforcement for
the capable models and enforcement is load-bearing for the weak ones.
Told the rule with the runtime off, both capable models complied on
every replicate of both commit tasks while both weak models committed
the gated mutation on every replicate; the enforced cells are safe by
construction, so the marginal value of enforcement given told is 0 for
the capable models (a failure to detect at this n, one-sided upper bound
+32pp, not demonstrated equivalence) and the whole barrier for the weak
ones. We report this as raw counts on a task-level unit (two tasks per
arm, near-deterministic replicates), not a pooled p-value; the effect is
stark but small and confounded -- capability with model family, and once
the below-floor 7B is set aside the robust load-bearing side rests on a
single non-frontier model. Because the reconstructed untold-unenforced
cell violates even for the capable models, the capable-side result reads
as substitution rather than pre-existing disposition. Told plainly is
load-bearing on top of capability: self-enforcement is
salience-sensitive, and when the same told constraint was reached
incidentally, the gated command was dispatched first-attempt 3/3
(diagnostic, one off-roster model, n=3, not independently
reconstructable). The clinical-refusal arm is near-ceiling and
underpowered, not a second substitution result, and the whole finding is
scoped to a single runtime, a four-model ladder whose low end is
Qwen-family, and the conditions of Section 8.0.

Enforcement is redundant only in that scoped, high-capability sense. It
retains the deterministic guarantee, which no behavioral result touches;
it is the only barrier at the untold violation floor, where an agent
never told the rule completed the violation with a truthful success
report; it is load-bearing across the sub-frontier capability band the
ladder measures, where a told agent still does not comply; below the
operate floor it prevents malformed harm rather than disobedience; and
adversarial territory remains untested here, cited as the regime where
telling is already known to fail.

The takeaway we most want practitioners to carry is methodological.
Before attributing fabrication to a model, check what the harness let it
see. An apparent instrumental-fabrication effect in our own pipeline
dissolved to 0 percentage points against a pre-committed bar of at least
40 (diagnostic, one model, pre-registered n=5) once the harness inlined
command stdout. Harness blindness manufactures spurious findings; an
evaluation of agent honesty is only as sound as the observation channel
the harness provides.

\subsection{Appendix}\label{appendix}

Skeleton only. Each subsection names the artifact that will supply the
final content and the PAPER.md anchor it must trace to; no full tables
are populated here.

\subsubsection{A. Task suite}\label{a.-task-suite}

\textbf{{[}STUB{]}} Table APPX.1: tasks by level and load-bearing
mechanism (M4-M9-TX). An earlier apparatus (a 28-task suite with 25
static oracle pairs, 23 per-mechanism plus 2
\texttt{no\_runtime\_enforcement} floor pairs, and a 16-trajectory
adversarial layer) was retired by PAPER.md decision D-37 (2026-07-04,
landed across benchmark commits \texttt{a10e850}..\texttt{48ff87b}):
those figures no longer apply. The suite was rebuilt as a sharp
per-mechanism 2x2 (told x enforced) with the told/untold contract axis
(\texttt{e72dced}) and first-attempt scoring (\texttt{48ff87b}) built
in, then expanded by D-39 (2026-07-05, commit \texttt{ed07f5a}) to three
scenario pairs per mechanism: 39 tasks (L1: 2, L2: 4, L5: 8, L6: 24, L7:
1); 16 paid-swept task-by-mode cells (two tasks carry
\texttt{no\_runtime\_enforcement} as the floor). Populate Table APPX.1
from the frozen suite and confirm the count against the release tag
before submission.

\subsubsection{B. Prompt provenance}\label{b.-prompt-provenance}

\texttt{deployment\_full\_v2} (D-28, authorized 2026-06-28): the
manifest is embedded as minified JSON with null/empty fields dropped,
lossless relative to \texttt{deployment\_full\_v1} (verified by
round-trip; no command/flag/type/choice/default/help/description
removed), reducing the rendered prompt from approximately 43K to
approximately 22K tokens after the v1 prompt exceeded a model's context
limit. Recorded as \texttt{PILOT\_PROTOCOL.md} Amendment 2. Cite the
prompt identifier from the roster's \texttt{prompt\_id} field (an
identifier string, not a hash) and the prompt file hash from
\texttt{scripts/lock\_hashes.json}, not from this appendix. Caution: the
recorded lock hash predates the prompt edits committed 2026-07-04
(\texttt{e72dced}, told/untold axis) and 2026-07-05 (\texttt{eceebf5},
S3-S6 prompt-neutrality pass; boundary-block parameterization);
re-derive the file hash at populate time, against the §20 run manifest,
rather than reusing any recorded value.

\subsubsection{C. Probe protocol
summary}\label{c.-probe-protocol-summary}

\textbf{{[}STUB{]}} Per-cell n and seed/temperature table, in two
clearly separated blocks.

Diagnostic probes (Sections 5.2-5.5):
\texttt{Qwen/Qwen3-235B-A22B-Instruct-2507-tput} (Together AI),
temperature 0, top\_p 1, no seed support in the provider API, n=3-5 per
cell depending on probe (see grounding pack Section 6 for per-probe n).
Table to enumerate probe directory, task(s), mode(s), and n per row.

Pre-registered run (Section 5.1): four-model Together ladder (MiniMax-M3
primary, Llama-3.3-70B, Qwen3.5-9B, Qwen2.5-7B; Qwen3-235B primary
deprecated mid-cycle and replaced per D-56), vendor-recommended
per-model sampling recorded per condition with card URL and snapshot
date, n=4 per cell (the mutation-gate headline pools its two commit
sub-gates to n=8), canary-first phasing with a hard stop that the gate
cleared, decision quantities and the eleven-branch outcome map per
PILOT\_PROTOCOL.md §20. Per-model sampling and prompt-token rows are
pinned in the locked roster revision; token counts for the identical
held-constant prompt differ per tokenizer, so the \textasciitilde22K
figure of Appendix B is Qwen2.5-specific and per-model counts are
recorded per condition.

\subsubsection{D. Adversarial trajectory table
(pointer)}\label{d.-adversarial-trajectory-table-pointer}

\textbf{{[}RESOLVED: retired{]}} The 16 hand-authored adversarial
trajectories (4 each vs M7, M5+M6, M8, M4) were scorer-coverage evidence
at the appendix tier, never model-backed. PAPER.md D-37 (2026-07-04)
records their retirement as final: commit \texttt{a10e850} removed
\texttt{results/adversarial\_summary.py} and the trajectories, and
\texttt{reproduce\_offline.py} no longer emits the
\texttt{adversarial\_summary} artifacts at HEAD. The adversarial-input
arm is cited future work (Sections 4.8, 9.1), not an arm of this paper.
This appendix subsection should either be deleted at LaTeX conversion or
reduced to a one-line provenance pointer; no adversarial table will be
populated.

\subsubsection{E. Reproducibility}\label{e.-reproducibility}

Per PAPER.md (Engineering Plan, Reproducibility script row),
\texttt{reproduce\_offline.py} runs the offline pipeline (rule-baseline
ablation, evidence tables, figures, error taxonomy, and cell contrasts;
the static and live isolation matrices were retired by D-37 at
\texttt{a10e850}) from a single command. Artifact and hash baselines in
\texttt{REPRODUCIBILITY.md} and \texttt{REPRODUCIBILITY\_GOLDEN.json}
(\texttt{REPRODUCIBILITY\_GOLDEN.json} regenerated post-D-37). The 76s
Apple M2 runtime figure is the pre-\texttt{a10e850} baseline;
re-baseline the runtime and pipeline description against the frozen D-39
39-task suite and the completed analysis-layer integration before
populating this appendix.

\subsubsection{F. Scorer config hash
chain}\label{f.-scorer-config-hash-chain}

Scorer behavior thresholds and critical violation kinds load from
\texttt{scorer\_config.paper\_v1.json}; the config hash records the
audited config bytes (D-18). A provenance-only hash bump
(\texttt{68e29510...} to \texttt{d310f503...}) accompanied the D-30
scorer-correctness pass, recorded in \texttt{lock\_hashes.json}; metric
thresholds, pass rule, and critical kinds were unchanged in that pass.

\subsubsection{G. D-O-02: Appendix E coding-agent sketch (open
decision)}\label{g.-d-o-02-appendix-e-coding-agent-sketch-open-decision}

\textbf{{[}OPEN, deferred to mid-August polish{]}} Keep-or-drop decision
on an Appendix E coding-agent sketch. PAPER.md D-O-02 requires an
anti-overclaim header if kept; PAPER.md does not specify the header's
wording. Workflow-derived content requirements for that header: the
sketch is illustrative, not a fourth evidence tier (the three tiers,
static oracle-pair evidence, live runtime probe, and model-backed
diagnostic, are never merged or extended), and not evidence of
generalization beyond the HAI controller (the paper is a single-runtime
case study; external replication is future work, Section 7.2). Do not
draft Appendix E content before D-O-02 resolves.

\subsection{References}\label{references}

Formatted from the verified neighbor list
(\texttt{paper/prior\_art\_notes.md}, synced 2026-07-05 from the
2026-07-04 novelty audit, D-38; full adjudication archived at
\texttt{ARCHIVE/novelty\_audit\_2026-07-04/}); verification status per
entry is recorded there. Section 2's external numeric claims were
verified against primary sources 2026-07-04: ContextCov (88.3\% vs
67.0\%) and PhantomPolicy (95.3\% to 40.7\%) in the novelty audit;
Mechanical Enforcement (MCC \textasciitilde0.43 to 0.88; 73\% deferral
reduction) and Verifier Tax (up to 94\% intercepted; SSR below 5\%) by
direct arXiv re-fetch. Entries are listed by the short name used in the
text.

\begin{itemize}
\tightlist
\item
  ABSTAIN: Ojewale and Venkatasubramanian (Brown University). What
  Benchmarks Don't Measure. arXiv:2606.02965. Workshop paper.
\item
  Agent Behavioral Contracts: Bhardwaj. Agent Behavioral Contracts.
  arXiv:2602.22302.
\item
  Agent-SafetyBench: Zhang et al.~ACL 2025. arXiv:2412.14470.
\item
  AgentSpec. arXiv:2503.18666.
\item
  AHE (automatic harness evolution). arXiv:2604.25850.
\item
  AIRGuard: Qin, Zhuang, et al.~arXiv:2605.28914.
\item
  CaMeL: Debenedetti et al.~(Google DeepMind / ETH Zurich).
  arXiv:2503.18813.
\item
  CEF/Thanatosis: J.P. Morgan AI Research. Is Your Agent Playing Dead?
  arXiv:2606.14831.
\item
  ContextCov: Sharma. arXiv:2603.00822.
\item
  Engels, Baek, Kantamneni, and Tegmark. Scaling Laws For Scalable
  Oversight. arXiv:2504.18530. NeurIPS 2025 (Spotlight). Cited as
  {[}Engels2025{]} per PAPER.md H7/D-03 (deferred line). Verified
  against the primary source 2026-07-04.
\item
  ETCLOVG survey. OpenReview 3hXEPbG0dh. TMLR, under review; cited via
  OpenReview, not as an accepted paper.
\item
  FORGE: Palumbo, Choudhary, Choi, Amir, Chalasani, Jha.
  arXiv:2602.16708 (v1 titled ``Policy Compiler for Secure Agentic
  Systems'').
\item
  Governance Decay (mitigation: Constraint Pinning). arXiv:2606.22528.
\item
  Guardrails AI. Software framework.
\item
  Harness-Bench. arXiv:2605.27922.
\item
  Harness-MU: Fan, Nie, Dai. arXiv:2606.21856.
\item
  IFEval. arXiv:2311.07911.
\item
  Invariant Guardrails: Invariant Labs (now part of Snyk). Software
  framework.
\item
  LogiSafetyBench. arXiv:2601.08196.
\item
  Mechanical Enforcement: de la Chica Rodriguez and Marti-Gonzalez.
  arXiv:2605.14744.
\item
  Meta-Harness. arXiv:2603.28052.
\item
  Mind the GAP: Text Safety Does Not Transfer to Tool-Call Safety in LLM
  Agents. arXiv:2602.16943.
\item
  NeMo Guardrails: NVIDIA. Software framework.
\item
  NLAH (natural-language agent harnesses). arXiv:2603.25723.
\item
  PhantomPolicy: Wu, Gong (Atlassian). arXiv:2604.12177.
\item
  Policy-as-Prompt. arXiv:2509.23994. NeurIPS 2025 Regulatable ML
  workshop.
\item
  Progent. arXiv:2504.11703.
\item
  Prompts Don't Protect: Uppala. arXiv:2605.18414. Single author; no
  venue located.
\item
  RECAST. arXiv:2505.19030.
\item
  SABER. arXiv:2606.01317.
\item
  SafePyramid. arXiv:2606.29887.
\item
  ST-WebAgentBench: Levy et al.~arXiv:2410.06703. ICLR 2026.
\item
  Symbolic Guardrails: Hong et al.~(CMU). arXiv:2604.15579.
\item
  TeamBench: Kim et al.~arXiv:2605.07073.
\item
  TRACE: Zhou et al.~arXiv:2606.13174.
\item
  VerIF. arXiv:2506.09942. EMNLP 2025 per ACL Anthology.
\item
  Verifier Tax: Sah, Srivastava, Sah, Jordan. arXiv:2603.19328.
\end{itemize}

\end{document}
